% methods.tex -- full text (2026-08-05)
% "A certified continuation machine for Mahler measure identities at CM points,
%  with seven new proofs of conjectures of Samart"
\documentclass[11pt]{article}
\usepackage{amsmath,amssymb,amsthm,mathtools}
\usepackage[margin=1in]{geometry}
\usepackage{booktabs}
\usepackage{hyperref}

% ---- notation (unified with formal.tex / formal2.tex) ----
\newcommand{\m}{\mathrm{m}}                 % logarithmic Mahler measure
\newcommand{\mt}{\widetilde{\mathrm{m}}}    % holomorphic (modified) Mahler measure
\newcommand{\Cc}{\mathbb{C}}
\newcommand{\Qq}{\mathbb{Q}}
\newcommand{\Zz}{\mathbb{Z}}
\newcommand{\Rr}{\mathbb{R}}
\newcommand{\Hh}{\mathbb{H}}
\newcommand{\Tt}{\mathbb{T}}
\newcommand{\Nn}{\mathbb{N}}
\newcommand{\ii}{\mathrm{i}}
\newcommand{\id}{\mathrm{d}}
\newcommand{\e}{\mathrm{e}}
\newcommand{\OO}{\mathcal{O}}
\newcommand{\cL}{\widetilde{L}}
\newcommand{\ReT}{\mathrm{Re}}
\newcommand{\ImT}{\mathrm{Im}}
\DeclareMathOperator{\EK}{EK}
\DeclareMathOperator{\Ast}{Ast}
\DeclareMathOperator{\dist}{dist}

\newtheorem{theorem}{Theorem}[section]
\newtheorem{lemma}[theorem]{Lemma}
\newtheorem{proposition}[theorem]{Proposition}
\newtheorem{corollary}[theorem]{Corollary}
\theoremstyle{definition}
\newtheorem{definition}[theorem]{Definition}
\newtheorem{example}[theorem]{Example}
\theoremstyle{remark}
\newtheorem{remark}[theorem]{Remark}

\title{A certified continuation machine for Mahler measure\\
identities at CM points, with seven new proofs of\\
conjectures of Samart}
\author{Huimin Zheng\thanks{%
College of Information and Network Engineering,
Anhui Science and Technology University,
Fengyang, Anhui 233100, P.~R.~China.
Email: \texttt{zhhm@ahstu.edu.cn}.}}
\date{}

\begin{document}
\maketitle

% ------------------------------------------------------------
\section*{Declaration on the use of AI tools}
The research reported in this article --- including the computational
exploration, the discovery of the proof strategy, the machine-certified
verifications, and the preparation of the manuscript --- was carried out
by the author with the assistance of the AI system \emph{Kimi}
(Moonshot AI). All mathematical content, including every proof and
every certified computation, has been checked and verified by the
author, who takes full responsibility for the correctness and integrity
of the article. All certification scripts are available for independent
verification (see Section~\ref{sec:cert}).
% ------------------------------------------------------------

\begin{abstract}
We axiomatize the differential-comparison continuation method of
\cite{Zf} into a general machine for proving Mahler measure identities
at CM points: a modular parametrization with rigorous tail bounds, a
single-valued holomorphic Mahler differential on the complement of the
critical image, a propagation lemma that replaces all monodromy and
universal-cover arguments, interval-arithmetic path certification, and
an exact three-track CM evaluation. Applied to Samart's family
$n_4(s)=4\m\bigl(x^4+y^4+z^4+1+s^{1/4}xyz\bigr)$, the machine proves
seven further entries of his 2015 table \cite[Table~6]{Sa15}: the
conjugate pairs at $s=8292456\pm3132675\sqrt7$ (whose two identities
are proved separately, not merely as a sum), $s=3656\pm2600\sqrt2$,
$s=-144$, $s=143208+101574\sqrt2$, and the value
$s=1207368+853632\sqrt2+697680\sqrt3+493272\sqrt6$ attached to the
class-number-two field $\Qq(\sqrt{-6})$ --- the first Mahler measure
identity in the $n_4$ family proved at a \emph{non-degenerate} CM
point of class number two (Samart's $s=2304$ \cite{Sa13}, attached
to $\tau=\sqrt{-6}/2$, is the degenerate value
$s_4(\ii\sqrt6/2)=s_4(\ii\sqrt6/6)$), where the lattice
sums involve two newforms and a genus character. New methodological
ingredients include: the Fricke ratio as the only non-vacuous
level/root-number discriminator (the ``functional-equation residue''
test is vacuous for root number $+1$ at every level); a genuine
twist-level trap ($g_{16}\otimes\chi_8$ has level $64$, not $32$, and
the wrong level silently scales the $L$-value); a genus-character
decomposition of the lattice sums in class number two; and the
discovery that Samart's applicability boundary $\{\ImT\tau=1/\sqrt2\}$
is not the topological boundary of the continuation region $V_4$.
Every identity is certified by an exact algebraic track (rational
arithmetic only, up to the quoted CM-theory inputs disclosed in
Section~\ref{sec:cert}), rigorous interval locks (half-widths down to
$9.1\times10^{-53}$), and independent $50$--$60$-digit numerical
cross-checks; all certification scripts are public. Finally we score
the ten remaining open entries of Table~6 against the machine's five
stations: five are within current reach (only mechanical work remains),
and five are blocked by genuine obstructions --- a two-dimensional
critical image or a wrong-sheet Eisenstein--Kronecker series ---
which we analyze precisely.
\end{abstract}

\medskip
\noindent\textbf{Keywords:} Mahler measure $\cdot$ modular forms
$\cdot$ complex multiplication $\cdot$ special values of
$L$-functions $\cdot$ rigorous interval arithmetic

\medskip
\noindent\textbf{Mathematics Subject Classification (2020):}
Primary 11R06; Secondary 11F03, 11M41, 11Y40

% ============================================================
\section{Introduction}\label{sec:intro}

\subsection{The question}

Let $\m(P)$ denote the logarithmic Mahler measure of a Laurent
polynomial and let
\[
  n_4(s):=4\,\m\bigl(x^4+y^4+z^4+1+s^{1/4}xyz\bigr),\qquad s\in\Cc,
\]
the value being independent of the choice of the fourth root
(Section~\ref{sec:machine}). In \cite{Sa13,Sa15} Samart expressed
$n_4(s)$, for $s$ in the image of a modular parametrization
$s_4(\tau)$, as an Eisenstein--Kronecker (EK) series, proved the
identity for pure imaginary $\tau$ with $\ImT\tau\ge1/\sqrt2$, and
conjectured on numerical evidence the twenty-nine $L$-value
evaluations of his Table~6 \cite{Sa15}: each attaches to a CM value
of $s$ a rational linear combination of $L$-values of CM weight-three
newforms and of quadratic Dirichlet characters.

The table is now largely, but not entirely, resolved. Eleven entries
were proved previously: the exterior values $s=256$ (Rogers
\cite{Ro}), $s=648,2304,20736,614656$ (Samart \cite{Sa13}),
$s=26856\pm15300\sqrt3$ (Samart's own Theorem~3.1 \cite{Sa15}), the
outside-the-critical-locus value $s=-3969$ (Fei \cite{Fei}), and
$s=-1024,-12288,-82944$ (Guo--Peng--Qin \cite{GPQ}). One entry is
\emph{refuted}: the conjecture $n_4(81)=40M_7$ fails numerically as a
statement about the true Mahler measure, the conjectured value being
the evaluation of the EK series at the attached CM point on a wrong
sheet \cite{Zf2}. He--Ye \cite{HY} resolved the companion family
$n_3$ completely; the $n_2$ family was completed in \cite{Zf,ZGQ,GPQ}.
Fei \cite{Fei} proved, for all rational singular moduli in twenty-three
families, identities at the level of the real part of the
\emph{holomorphic} Mahler measure $\ReT\mt$; at parameters inside the
critical locus the passage from $\ReT\mt$ to the genuine Mahler
measure is not automatic --- a point stressed already by Rodriguez
Villegas \cite{RV} --- and it is precisely this passage that the
continuation machine of \cite{Zf} and of the present paper performs.

\begin{remark}[The status of $s=-3969$]\label{rem:fei}
For the entry $s=-3969$ above, \cite{Fei} proves the identity for
$\ReT\mt$ only, and the determination of the region off which
$\m=\ReT\mt$ is left open there (his Remark~2.4); by
Section~\ref{sec:boundary} below, avoiding the critical image alone
does not suffice (the wrong-sheet phenomenon). The missing step is
supplied by the present machine: the attached point
$\tau=(1+\sqrt{-7})/2$ lies on the vertical segment
$\ReT\tau=1/2$, $0.702\le\ImT\tau\le\sqrt{21}/2$, which
\texttt{n5\_line\_cert.py} certifies inside the good component $W$
by interval arithmetic (public repository; the certificate was
produced for the Phase-8 entries of Table~\ref{tab:open}). Hence
Theorem~\ref{thm:main} applies and the $s=-3969$ evaluation holds
for the genuine Mahler measure.
\end{remark}

The purpose of this paper is threefold:
\begin{enumerate}
\item[(i)] to axiomatize the differential-comparison continuation
  method of \cite{Zf} into a general, family-independent machine
  (Section~\ref{sec:machine}), with an explicit certifier architecture
  whose trusted base consists of pure mathematics, documented
  interval-arithmetic semantics, and public source code
  (Section~\ref{sec:cert});
\item[(ii)] to apply the machine to seven further entries of Samart's
  Table~6 (Theorems A--F below), covering all four regimes that occur:
  deep pure-imaginary interior points (A, C, D, E), a non-pure-imaginary
  interior point (B), the first class-number-two CM point (E), and a
  point on Samart's applicability boundary that turns out to be a strict
  interior point of the continuation region (F);
\item[(iii)] to take stock: to score the ten remaining open entries of
  Table~6 against the machine's checklist (\S\ref{subsec:inventory})
  and to describe precisely where and why the method stops
  (Section~\ref{sec:boundary}).
\end{enumerate}

\subsection{Results of this paper}

Throughout, $d_k:=L'(\chi_{-k},-1)$ for the odd quadratic Dirichlet
character $\chi_{-k}$ of conductor $k$, and $M_k$ or $M_k^{(i)}$
denotes $L'(g_k,0)$, resp.\ $L'(g_k^{(i)},0)$, for the indicated
weight-$3$ CM newform(s); a superscript $\mathrm{tw}$ denotes the
$L'$-value of an explicitly identified quadratic twist.

\begin{theorem}[Theorem A: the $\sqrt7$ pair]\label{thm:A}
Let $M_7:=L'(g_7,0)$ with
$g_7(\tau)=\eta(\tau)^3\eta(7\tau)^3$ \textup{(}LMFDB
\textup{\textsf{7.3.b.a}}\textup{)} and
$M_7^{\mathrm{tw}}:=L'(g_7\otimes\chi_{-4},0)$. Then
\begin{align}
n_4\bigl(8292456+3132675\sqrt7\bigr)
  &=\frac{5}{28}\bigl(4M_7^{\mathrm{tw}}+224M_7+32d_4+7d_7\bigr),
    \label{eq:A1}\\
n_4\bigl(8292456-3132675\sqrt7\bigr)
  &=\frac{5}{14}\bigl(4M_7^{\mathrm{tw}}-224M_7+32d_4-7d_7\bigr).
    \label{eq:A2}
\end{align}
\end{theorem}

\begin{theorem}[Theorem B]\label{thm:B}
With $M_{12}:=L'(g_{12},0)$, $g_{12}(\tau)=\eta(2\tau)^3\eta(6\tau)^3$
\textup{(}LMFDB \textup{\textsf{12.3.c.a}}\textup{)},
\[
  n_4(-144)=\frac{10}{3}\bigl(4M_{12}+d_3\bigr).
\]
\end{theorem}

\begin{theorem}[Theorem C]\label{thm:C}
With $M_8:=L'(g_8,0)$,
$g_8(\tau)=\eta(\tau)^2\eta(2\tau)\eta(4\tau)\eta(8\tau)^2
\in S_3(\Gamma_0(8),\chi_{-8})$,
and $M_8^{\mathrm{tw}}:=L'(g_8\otimes\chi_8,0)$, where
$\chi_8=(2/\cdot)$ is the even primitive character of conductor $8$,
\[
  n_4\bigl(3656+2600\sqrt2\bigr)
  =\frac58\bigl(4M_8^{\mathrm{tw}}+28M_8+4d_4+d_8\bigr).
\]
\end{theorem}

\begin{theorem}[Theorem D]\label{thm:D}
With $M_{16}:=L'(g_{16},0)$,
$g_{16}(\tau)=\eta(4\tau)^6\in S_3(\Gamma_0(16),\chi_{-4})$, and
$M_{16}^{\mathrm{tw}}:=L'(g_{16}\otimes\chi_8,0)$
\textup{(}level $64$ --- see \S\ref{subsec:twisttrap}\textup{)},
\[
  n_4\bigl(143208+101574\sqrt2\bigr)
  =\frac{5}{16}\bigl(4M_{16}^{\mathrm{tw}}+20M_{16}+9d_4+4d_8\bigr).
\]
\end{theorem}

\begin{theorem}[Theorem E: class number two]\label{thm:E}
Let $K=\Qq(\sqrt{-6})$ \textup{(}class number $2$\textup{)}, and let
$g_{24}^{(1)},g_{24}^{(2)}$ be the two newforms of
$S_3(\Gamma_0(24),\chi_{-24})$, theta series of the two
gr\"ossencharacters of $K$ of type $(2,0)$ \textup{(}LMFDB
\textup{\textsf{24.3.h.a}} and \textup{\textsf{24.3.h.b}}\textup{)}.
With $M_{24}^{(i)}:=L'(g_{24}^{(i)},0)$ and
$M_{24}^{(i),\mathrm{tw}}:=L'(g_{24}^{(i)}\otimes\chi_{-8},0)$
\textup{(}level $96$\textup{)},
\begin{multline*}
  n_4\bigl(1207368+853632\sqrt2+697680\sqrt3+493272\sqrt6\bigr)\\
  =\frac{5}{48}\Bigl(4M_{24}^{(1),\mathrm{tw}}
   +4M_{24}^{(2),\mathrm{tw}}+28M_{24}^{(1)}+12M_{24}^{(2)}
   +28d_3+24d_4+8d_8+d_{24}\Bigr).
\end{multline*}
\end{theorem}

\begin{theorem}[Theorem F: a boundary point that is interior]\label{thm:F}
With the notation of Theorem~\ref{thm:C},
\[
  n_4\bigl(3656-2600\sqrt2\bigr)
  =\frac54\bigl(4M_8^{\mathrm{tw}}-28M_8+4d_4-d_8\bigr).
\]
The $s$-value is the conjugate of Theorem~\ref{thm:C}, but the attached
CM point $\tau_4=(1+\sqrt{-2})/2$ lies \emph{exactly} on Samart's
applicability boundary $\ImT\tau=1/\sqrt2$, where his proof does not
apply; the identity is proved by the continuation machine
\textup{(}Section~\ref{sec:app5}\textup{)}, which shows that $\tau_4$
is a strict interior point of the continuation region $V_4$
\textup{(}certified margin $0.11590$\textup{)}.
\end{theorem}

\begin{theorem}[The machine, informal]\label{thm:machine}
\textup{(}Precise form: Theorem~\ref{thm:main}.\textup{)}
Let $P_c$ be a tempered Laurent-polynomial family with a modular
parametrization $c=c(\tau)$ and critical image $K\subset\Cc$, and let
$E(\tau)$ be the Eisenstein--Kronecker expression for the holomorphic
Mahler measure, valid on an open set $U\subset\Hh$. If
\textup{(i)} the period integral $\Omega(c)$ is single-valued and
holomorphic on $\Cc\setminus K$, \textup{(ii)} $c(\tau_0)$ is a CM
parameter with $\tau_0\in\overline W$, where $W$ is the connected
component of $\Hh\setminus c^{-1}(K)$ containing $U$, and
\textup{(iii)} the lattice sums underlying $E(\tau_0)$ evaluate
exactly, then $\m(P_{c(\tau_0)})$ equals the corresponding
$L$-value expression; the inclusion $\tau_0\in\overline W$ is
certifiable by rigorous interval arithmetic.
\end{theorem}

All seven identities are machine-certified. Concretely, each proof
consists of an exact algebraic track (rational/algebraic-integer
arithmetic only, no floating point), rigorous interval locks of the
CM values and of the final identities (half-widths between
$9.1\times10^{-53}$ and $2.3\times10^{-51}$), and independent
$50$--$60$-digit numerical cross-checks; the scripts are listed in
Section~\ref{sec:cert} and are publicly available
(see the Data Availability statement). Two representative reference
values:
\begin{align*}
  M_8&=0.1353184954231269106150060654595395883532\ldots,\\
  M_8^{\mathrm{tw}}
   &=1.549084847095611486271658159444376281785\ldots
     \quad(\text{level }32,\ w=+1),
\end{align*}
computed from the functional equations and cross-checked against
direct lattice-sum evaluations; further values are collected in the
application sections.

\subsection{The ten remaining entries}\label{subsec:inventory}

Table~6 of \cite{Sa15} has twenty-nine entries: eighteen are now
proved (eleven previously known, plus the seven of Theorems
A--F), one is refuted ($s=81$ \cite{Zf2}), and ten remain open.
The machine reduces a proof of each entry to five \emph{stations}
(Section~\ref{sec:machine} for definitions):
\begin{enumerate}
\item[\textbf{S1}] \emph{Parametrization}: a CM point $\tau_0$ with
  $s_4(\tau_0)=s$ identified, at least numerically;
\item[\textbf{S2}] \emph{Critical image}: the position of
  $c(\tau_0)$ relative to the astroid disc $\Ast$ and the cross
  $[-4,4]\cup\ii[-4,4]$;
\item[\textbf{S3}] \emph{Interior-or-boundary}: membership of
  $\tau_0$ in $\overline W$, certified by a \texttt{cert2} path,
  or direct applicability of Samart's theorem;
\item[\textbf{S4}] \emph{\texttt{cert0} algebraicity}: an exact
  algebraic lock of the value $s_4(\tau_0)$;
\item[\textbf{S5}] \emph{P1 assemblability}: an exact CM evaluation
  of $\EK_4(\tau_0)$ in the relevant imaginary quadratic field.
\end{enumerate}
Table~\ref{tab:open} scores the ten open entries against these
stations.

\begin{table}[htb]
\centering\scriptsize
\caption{The ten open entries of Samart's Table~6, scored against the
five stations of the machine. ``Outside/inside'' refers to the
astroid disc $\Ast$; $\zeta$ abbreviates
$1207368$, $a=853632$, $b=697680$, $d=493272$ in rows 2--4:
$s=\zeta\pm a\sqrt2\pm b\sqrt3\pm d\sqrt6$ with the indicated signs.}
\label{tab:open}
\setlength{\tabcolsep}{2.5pt}
\begin{tabular}{cp{2.3cm}llp{1.9cm}p{2.3cm}p{2.5cm}l}
\toprule
\# & $s$ & $\tau_0$ & $\ImT\tau_0$ & S2: $c(\tau_0)$ & S3: position
 & S4/S5 status & verdict\\
\midrule
1 & $143208-101574\sqrt2$ & $(1+2\ii)/2$ & $1$
  & outside $\Ast$ ($|s|>256$)
  & interior; $\tau_0$ lies on the certified path of
    Section~\ref{sec:app5}
  & S4 done (conjugate of Thm.~D lock); S5 ready ($\Qq(\ii)$, $h=1$)
  & ready\\
2 & $\zeta$, signs $({+}{-}{-})$ & $(1+\ii\sqrt6)/2$ & $1.225$
  & outside $\Ast$ ($|s|>256$)
  & interior; short \texttt{cert2} extension
  & S4 done (\texttt{cert0\_n4\_p4\_t3.py}); S5 ready ($h=2$,
    \S\ref{sec:app4})
  & ready\\
3 & $\zeta$, signs $({-}{+}{-})$ & $\ii\sqrt6/4$ & $0.612$
  & outside $\Ast$ ($s>256$)
  & \emph{exterior}: EK series wrong-sheeted below $1/\sqrt2$
  & S4 done (same lock); S5 moot
  & blocked\\
4 & $\zeta$, signs $({-}{-}{+})$ & $(2+\ii\sqrt6)/4$ & $0.612$
  & \emph{inside} $\Ast$ ($s\approx-2.45$)
  & blocked (two-dimensional image, open mapping)
  & S4 done (same lock); S5 moot
  & blocked\\
5 & $(-192303+85995\sqrt5)/2$ & $(3+\ii\sqrt{15})/6$ & $0.646$
  & inside $\Ast$ ($s\approx-6.17$)
  & blocked (two-dimensional image); also exterior
  & S4 open (disc.\ $-15$, $h=2$); S5 needs level-$15$ pair
  & blocked\\
6 & $(-192303-85995\sqrt5)/2$ & $(1+\ii\sqrt{15})/2$ & $1.936$
  & outside $\Ast$
  & interior; \texttt{cert2} feasible on $\ReT\tau=1/2$
  & S4 open (pair lock with \#5); S5 ready ($h=2$ machinery)
  & ready\\
7 & $-893952+516096\sqrt3$ & $(3+\ii\sqrt{21})/6$ & $0.764$
  & outside $\Ast$ (functional $\approx3.02$)
  & interior; $\tau_0$ lies on the certified path of
    Section~\ref{sec:app5}
  & S4 open (pair lock with \#8); S5 ready (level $84$, $h=2$)
  & ready\\
8 & $-893952-516096\sqrt3$ & $(1+\ii\sqrt{21})/2$ & $2.291$
  & outside $\Ast$
  & interior; \texttt{cert2} feasible
  & S4 open (same pair lock); S5 ready
  & ready\\
9 & $347648256+141926400\sqrt6$ & $\ii\sqrt{42}/42$ & $0.154$
  & outside $\Ast$ ($c>4$ real)
  & \emph{exterior}: wrong sheet, deep below $1/\sqrt2$
  & S4 open (disc.\ $-168$, $h=4$); S5 moot
  & blocked\\
10 & $347648256-141926400\sqrt6$ & $\ii\sqrt{42}/14$ & $0.463$
  & outside $\Ast$ ($s\approx995$, $c>4$ real)
  & \emph{exterior}: wrong sheet below $1/\sqrt2$
  & S4 open; S5 moot
  & blocked\\
\bottomrule
\end{tabular}
\end{table}

The conjectured right-hand sides are Samart's \cite[Table~6]{Sa15};
for example, entry \#1 is
$n_4(143208-101574\sqrt2)\stackrel?=\frac58\bigl(4M_{16}^{\mathrm{tw}}
-20M_{16}-9d_4+4d_8\bigr)$, the conjugate sign pattern of
Theorem~\ref{thm:D}, and entries \#2--\#4 are the three remaining
sign patterns of Theorem~\ref{thm:E}'s family. Five entries are
\emph{ready}: for \#1 and \#7 the target CM point even lies on the
already certified path of Section~\ref{sec:app5} (legs
$\{x=1/2,\ 0.702\le y\le1\}$ and $\{y=1,\ 0\le x\le1/2\}$), so only
stations S4/S5 remain; \#2, \#6, \#8 require a routine
\texttt{cert2} adaptation and a new \texttt{cert0} pair lock. The
five \emph{blocked} entries split into two genuine obstructions,
analyzed in Section~\ref{sec:boundary}: the two-dimensional critical
image (\#4, \#5: the $s=81$ mechanism of \cite{Zf2}), and the
wrong-sheet phenomenon of the EK series below $\ImT\tau=1/\sqrt2$
(\#3, \#9, \#10, and again \#5), where the conjectured identity
concerns the true Mahler measure but the series side of the machine
has no valid input.

Beyond Table~6, three further open directions are discussed in
Section~\ref{sec:boundary}: Samart's ``three modular $L$-value''
hypothesis $s=16+1600\sqrt[3]2-1280\sqrt[3]4$ \cite{Sa15}, which has
no attached CM point in the $s_4$-image; the exterior entries, for
which a corrected formula of Samart--Tao $\widetilde n$-type \cite{ST}
would be needed; and the Boyd-lineage two-variable conjectures, where
no CM evaluation is available at all.

% ============================================================
\section{The machine}\label{sec:machine}

This section axiomatizes the method of \cite[\S3--\S5]{Zf} in
family-independent language. The running example is the $n_4$-family,
but the only family-specific inputs are the parametrization, the
critical image, and the anchor identity; we display them as separate
\emph{stations} S1--S5, matching the checklist of
\S\ref{subsec:inventory}.

\subsection{Families, parametrizations, critical images}\label{subsec:fam}

\begin{definition}[Tempered family]\label{def:tempered}
A \emph{tempered family} is a Laurent-polynomial family
$P_c(x_1,\dots,x_n)=P(x_1,\dots,x_n)+c$, $c\in\Cc$, such that
$c\mapsto\m(P_c)$ is continuous on all of $\Cc$ --- for the families
here this holds by the elementary dominated-convergence argument of
\cite[Lemma~2.1]{Zf}. A \emph{critical image} is any closed set $K\subset\Cc$ containing
the full image $P(\Tt^n)$; on $\Cc\setminus K$ the polynomial $P_c$
does not vanish on $\Tt^n$, the integrand $\log|P_c|$ is
real-analytic in $c$, and $c\mapsto\m(P_c)$ is real-analytic.
(Taking only the critical \emph{values} of $P$ would not suffice:
$P_c$ vanishes on $\Tt^n$ for every $-c\in P(\Tt^n)$, including
regular values --- e.g.\ $P=x+x^{-1}$ on $\Tt^1$ has critical values
$\pm2$ but image $[-2,2]$. In both running examples below
$K=\overline{P(\Tt^n)}$, the closure of the full image.)
\end{definition}

Two families serve as running examples.

\begin{example}[The $n_2$ family \cite{Zf}]\label{ex:n2}
$P+c$ with $P=(x+x^{-1})(y+y^{-1})(z+z^{-1})$ and
$c=\sqrt{s_2(\tau)}$, $s_2(\tau)=q^{-1}\prod_{n\ge1}(1+q^{2n-1})^{24}$.
The critical image in the $c$-plane is the slit $[-8,8]$, and in the
$s$-plane the segment $[0,64]$: one-dimensional.
\end{example}

\begin{example}[The $n_4$ family]\label{ex:n4}
$P_c=(x^4+y^4+z^4+1)/(xyz)+c$ after clearing the monomial (Mahler
measure is unchanged), with the natural \emph{holomorphic}
parametrization
\begin{equation}\label{eq:cdef}
  c(\tau)=\Bigl(\frac{\eta(2\tau)}{\eta(\tau)}\Bigr)^{6}
   \bigl(16A(\tau)^4+A(\tau)^{-4}\bigr),\qquad
  A(\tau)=\frac{\eta(\tau)\eta(4\tau)^2}{\eta(2\tau)^3},\qquad
  s_4(\tau)=c(\tau)^4 .
\end{equation}
Both $c$ and $s_4$ are holomorphic on $\Hh$ --- no branch choice is
needed, in contrast to Example~\ref{ex:n2} --- and $s_4$ has an
integral $q$-expansion
$s_4=q^{-1}\bigl(1+104q+4372q^2+\cdots\bigr)\in\Zz(\!(q)\!)$
(exact integer check to order $q^{30}$ in
\texttt{cert0\_s4\_m144.py}). The critical image is
\emph{two-dimensional}: $P(\Tt^3)$ is the closed astroid disc
\begin{equation}\label{eq:astroid}
  \Ast=\bigl\{c\in\Cc:\ |\ReT c|^{2/3}+|\ImT c|^{2/3}\le 4^{2/3}\bigr\},
\end{equation}
since $P$ is a sum of four unit phasors with product $1$
\cite{Zf2}. The branch cut of the holomorphic Mahler measure is the
cross $\{c:c^4\in(0,256]\}=\bigl([-4,4]\cup\ii[-4,4]\bigr)
\setminus\{0\}$, contained in $\overline{\Ast}$; hence avoiding
$\Ast$ automatically avoids every bad set, and we take $K=\Ast$.
\end{example}

\begin{remark}\label{rem:2D}
The dimension of the critical image is the single most important
topological invariant of a family for this method. A one-dimensional
slit does not separate the plane: every point of the slit lies in the
closure of the good domain, and the propagation lemma below reaches
boundary points by continuity (this is what makes the $k=1$ proof of
\cite{Zf} possible). The two-dimensional astroid disc has interior:
a CM parameter landing inside it (\#4 and \#5 of
Table~\ref{tab:open}; the refuted $s=81$) lies inside an open bad
region, no path can approach it from the good domain (open mapping),
and the machine stops. See Section~\ref{sec:boundary}.
\end{remark}

\begin{definition}[The good domain and the EK expression]
For a holomorphic parametrization $c(\tau)$ set
\[
  V:=\bigl\{\tau\in\Hh:\ c(\tau)\notin K\bigr\}\qquad
  (\text{open, since }c\text{ is holomorphic}),
\]
and let $F(\tau):=\m\bigl(P_{c(\tau)}\bigr)$ (times the family's
normalization factor), continuous on all of $\Hh$ and real-analytic
on $V$. An \emph{EK datum} is a holomorphic function $E(\tau)$ on
$\Hh$ with $G:=\ReT E$ (in alternative normalizations $2\ReT E$ or
$\ImT E$; all that matters is that $G$ be the real part of a
holomorphic function) real-analytic everywhere, such that $F=G$ on a
nonempty open \emph{anchor} set $U\subset V$. For the $n_4$ family
(Example~\ref{ex:n4}), with $U_j$ as in \eqref{eq:Uj} below,
\begin{equation}\label{eq:E4}
  E_4(\tau):=-\ii\Bigl(2\pi\tau+\frac{10}{\pi^3}
   \bigl(U_1(\tau)-2U_2(\tau)\bigr)\Bigr),\qquad
  \EK_4(\tau):=\ReT E_4(\tau)
   =\ImT\Bigl[2\pi\tau+\frac{10}{\pi^3}\bigl(U_1-2U_2\bigr)\Bigr],
\end{equation}
holomorphic, resp.\ real-analytic, on all of $\Hh$; the double series
\begin{equation}\label{eq:Uj}
  U_j(\tau):=\sum_{m\neq0}\frac1m\sum_{n\in\Zz}(jm\tau+n)^{-3}
\end{equation}
converges absolutely, locally uniformly on $\Hh$
\cite[\S2.4]{Zf}. The lattice-sum form used for all CM evaluations is
\begin{equation}\label{eq:EK4lattice}
  \EK_4(\tau)=\frac{10\,\ImT\tau}{\pi^3}
   \bigl(-T_1(\tau)+4T_2(\tau)\bigr),\qquad
  T_d(\tau):={\sum_{\lambda\in\Lambda_d(\tau)}}^{\!\prime}
   \Bigl[\frac{2\ReT\bar\lambda^2}{|\lambda|^6}
        +\frac{1}{|\lambda|^4}\Bigr],
\end{equation}
with $\Lambda_d(\tau):=\Zz+\Zz\,d\tau$; the series is absolutely
convergent, and \eqref{eq:EK4lattice} follows from
$w^{-3}-\bar w^{-3}=-2\ii(3u^2v-v^3)|w|^{-6}$ and
$3u^2-v^2=4u^2-|w|^2$ exactly as in \cite[Lemma~2.6]{Zf}.
\end{definition}

\subsection{The Mahler differential is single-valued}\label{subsec:diff}

\begin{proposition}[Mahler differential formula]\label{prop:mahlerdiff}
Let $P_c$ be a tempered family and $K\subset\Cc$ its critical image.
Then
\[
  \Omega(c):=\int_{\Tt^n}\frac{\id\mu}{P_c}
\]
is single-valued and holomorphic on $\Cc\setminus K$, and as real
$1$-forms there,
\begin{equation}\label{eq:dm}
  \id\m(P_c)=\ReT\bigl[\Omega(c)\,\id c\bigr].
\end{equation}
\end{proposition}
\begin{proof}
For any compact $K_0\Subset\Cc\setminus K$ one has
$|P_c|\ge\dist(K_0,K)>0$ uniformly for $c\in K_0$ on $\Tt^n$, so
differentiation under the integral sign is legitimate and gives
$\Omega'(c)=-\int_{\Tt^n}\id\mu/P_c^2$; holomorphy and
single-valuedness follow (the integral is defined over the
\emph{fixed} cycle $\Tt^n$). Writing $c=c_1+\ii c_2$ and
differentiating $\m(P_c)=\int\log|P_c|\id\mu$ under the integral
sign gives
$\partial_{c_1}\m=\int\ReT(1/P_c)\id\mu=\ReT\Omega$ and
$\partial_{c_2}\m=\int\ReT(\ii/P_c)\id\mu=-\ImT\Omega$, hence
$\id\m=\ReT\Omega\,\id c_1-\ImT\Omega\,\id c_2
=\ReT[\Omega\,\id c]$. The sign is $+$ under the present $P+c$
convention (it is $-$ in the $f-c$ convention of
\cite[Lemma~2.7]{Zf}).
\end{proof}

\begin{remark}\label{rem:monodromy}
Proposition~\ref{prop:mahlerdiff} is the point at which all
monodromy and universal-cover arguments disappear from the method.
The holomorphic Mahler measure $\mt$ of Villegas \cite{RV} is
multivalued, but its multivaluedness lives in locally constant
imaginary periods: every branch satisfies $\mt'(c)=\Omega(c)$ on
$\Cc\setminus K$. Comparing \emph{differentials} rather than
functions therefore never sees a branch choice, and the identity
theorem below applies to a genuine single-valued holomorphic
function.
\end{remark}

\subsection{The propagation theorem}\label{subsec:prop}

\begin{lemma}[Propagation lemma]\label{lem:prop}
Let $U\subset\Cc$ be connected open, $Z\subset U$ relatively closed,
and $H:U\to\Rr$ continuous and locally constant on $U\setminus Z$.
Then $H$ is constant on every connected component $W$ of
$U\setminus Z$; if $H|_W\equiv C$ then
$H|_{\overline W\cap U}\equiv C$.
\end{lemma}
\begin{proof}
A locally constant function is constant on each connected component.
For the second assertion, let $p\in\overline{W}\cap U$ and choose
$p_n\in W$ with $p_n\to p$; continuity gives $H(p)=\lim_n H(p_n)=C$.
\end{proof}

\begin{remark}\label{rem:devil}
The conclusion cannot be upgraded to ``$H$ constant on $U$'': the
devil's staircase is continuous, locally constant off the Cantor set,
yet non-constant. Membership of the target point in $\overline W$ is
a genuine topological input --- supplied in our applications either
by Samart's proved region (Theorems A, C, D, E), or by an explicitly
certified path (Theorems B, F), and \emph{absent} for the blocked
entries of Table~\ref{tab:open}.
\end{remark}

\begin{theorem}[Main theorem: differential-comparison continuation]
\label{thm:main}
Let $P_c$ be a tempered family with critical image $K$, let
$c(\tau)$ be holomorphic on $\Hh$, set
$V=\{\tau:c(\tau)\notin K\}$, and let $F(\tau)=\m(P_{c(\tau)})$
\textup{(}with the family normalization\textup{)}. Let $E(\tau)$ be
holomorphic on $\Hh$, set $G:=\ReT E$ (the $n_4$ normalization
$G=\EK_4=\ReT E_4$ of \eqref{eq:E4}; rescaling $E$ absorbs any
other convention), and suppose $F=G$ on a nonempty
open subset of a connected component $W$ of $V$. Then
\begin{equation}\label{eq:mainconcl}
  F(\tau)=G(\tau)\qquad\text{for every }
  \tau\in\overline{W}\cap\Hh .
\end{equation}
\end{theorem}
\begin{proof}
Write $F=\nu\,\m(P_{c(\tau)})$ with the family's constant
normalization factor $\nu$ ($\nu=4$ for the $n_4$ family).
On $V$ both $\Omega(c(\tau))c'(\tau)$ and $E'(\tau)$ are holomorphic,
so $\Psi(\tau):=\nu\,\Omega(c(\tau))c'(\tau)-E'(\tau)$ is
holomorphic on $V$. By the chain rule applied to \eqref{eq:dm},
$\id F=\nu\,\ReT[\Omega(c(\tau))c'(\tau)\,\id\tau]$ on $V$, and
$\id G=\ReT[E'(\tau)\,\id\tau]$ there; hence
$\id(F-G)=\ReT[\Psi(\tau)\,\id\tau]$. On the anchor open set
$F-G=0$, so $\ReT[\Psi\,\id\tau]=0$ in all directions
$\id\tau$; testing $\id\tau=1$ and $\id\tau=\ii$ gives $\Psi=0$
there, and the identity theorem on the connected open set $W$ gives
$\Psi\equiv0$ on $W$. Thus $H:=F-G$ is locally constant on $W$,
constant since $W$ is connected, and zero because it vanishes on the
anchor; $H$ is continuous on all of $\Hh$ (Mahler continuity,
Definition~\ref{def:tempered}; $G$ is real-analytic everywhere), so
Lemma~\ref{lem:prop} with $U=\Hh$, $Z=c^{-1}(K)$ gives
$H\equiv0$ on $\overline W\cap\Hh$.
\end{proof}

\begin{remark}[The five stations]\label{rem:stations}
Instantiating Theorem~\ref{thm:main} for a conjectured identity
$n_4(s)=R$ at a CM parameter splits into the five stations of
\S\ref{subsec:inventory}: S1 identifies $\tau_0$ with
$s_4(\tau_0)=s$; S2 locates $c(\tau_0)$ relative to $\Ast$; S3
certifies $\tau_0\in\overline W$ (or gets it for free from Samart's
proved region); S4 upgrades the numerical coincidence $s_4(\tau_0)=s$
to an exact algebraic lock (\S\ref{subsec:cert0}); S5 evaluates
$\EK_4(\tau_0)$ exactly as a rational linear combination of
$L$-values (\S\ref{subsec:verify}). Theorem~\ref{thm:main} then
yields $n_4(s_4(\tau_0))=\EK_4(\tau_0)$, and S4+S5 turn this into
the identity $n_4(s)=R$.
\end{remark}

\subsection{The anchor identity for the \texorpdfstring{$n_4$}{n4}
family}\label{subsec:anchor}

The anchor input is Samart's theorem, quoted with one repair.

\begin{theorem}[Samart, {\cite[Prop.~2.1(iii)]{Sa13}}]\label{thm:samart4}
For pure imaginary $\tau=\ii y$ with $y\ge1/\sqrt2$,
$n_4\bigl(s_4(\tau)\bigr)=\EK_4(\tau)$.
\end{theorem}

\begin{lemma}[Composite continuation lemma: the anchor on
$|s_4|>256$]\label{lem:anchor}
Let $D\subset\Hh$ be open with $|s_4(\tau)|>256$ for all $\tau\in D$.
Then $n_4(s_4(\tau))=\EK_4(\tau)$ on $D$.
\end{lemma}
\begin{proof}[Proof sketch with quoted inputs]
Three quoted ingredients compose. (i) \emph{Rogers' evaluation}
(Prop.~2.2 and Thm.~2.3 of \cite{Ro}): the ${}_5F_4$ hypergeometric
expression for $n_4$ converges for $|s|>256$
($\sum b-\sum a=3/2>0$) and equals $n_4(s)$ there: both sides are
real parts of holomorphic functions of $s$ on $\Cc\setminus\Ast$,
equal on the real ray $s>256$ by Rogers' theorem, hence equal on the
connected open set $\{|s|>256\}$ by the identity theorem. (ii)
\emph{Rogers' transformation} \cite[Thm.~2.3]{Ro}:
$f_4(s_4(q))=-\tfrac13 G(q)+\tfrac23 G(q^2)$ is an identity of
functions analytic near $q=0$, extending to $|q|<1$ by the identity
theorem (the composition with (i) is then restricted to the
connected component of $\{\tau:|s_4(\tau)|>256\}$ containing
$\ii\infty$, which is where every application lies). (iii) \emph{Samart's computation} \cite[proof of
Prop.~2.1(iii)]{Sa13}: the identification of the $q$-series in (ii)
with the $U$-series expression $\EK_4$ uses termwise integration and
Fourier expansions whose convergence requires only $|q|<1$. The
single step of Samart's proof that genuinely needs real $q$ --- the
comparison $|s_4(q)|\ge s_4(\e^{-\pi\sqrt2})=256$ by real-variable
monotonicity, which locates his argument in the convergence region
--- is replaced here by the hypothesis $|s_4|>256$ on $D$, certified
by interval arithmetic in each application
(Theorems~\ref{thm:B} and~\ref{thm:F}). This repair was identified
by a line-by-line check of the proof of \cite[Prop.~2.1(i)]{Sa13}
and of the convergence analysis of \cite[Thm.~2.3]{Ro} (Samart's
proof of Prop.~2.1(iii) itself is one sentence, referring to the
(i)-proof); no other step uses the pure-imaginary hypothesis.
\end{proof}

Lemma~\ref{lem:anchor} upgrades Samart's theorem from a statement on
a half-line to an anchor on an open set: the anchor box
$D=\{|\ReT\tau|\le1/32,\ |\ImT\tau-1|\le1/32\}$ carries the
certified bound $|s_4|\ge493$ (Section~\ref{sec:app2}), so $F=G$ on
$D$, and Theorem~\ref{thm:main} propagates the identity through the
whole component $W$ of $V_4:=\{\tau:c(\tau)\notin\Ast\}$ containing
$D$. For the deep pure-imaginary points of Theorems A, C, D, E even
this is unnecessary: Theorem~\ref{thm:samart4} applies directly.

\subsection{The certifier: \texttt{cert0}, \texttt{cert2},
\texttt{verify\_P1}}\label{sec:cert}

The non-elementary computations are certified by interval arithmetic.
The object entering the proofs is not a program but a
\emph{certificate}: finite data (parameter blocks with strict
enclosures, or exact rational equalities) whose validity is checked
by deterministic, finite, independently re-runnable computations.
The trusted base consists of exactly three items:
(i) the mathematical lemmas below (tail bounds, inclusion property);
(ii) the documented outward-rounding semantics of the interval
arithmetic employed (\texttt{mpmath.iv} rounds interval endpoints
outward, with guard digits for the complex transcendental functions
\cite{mpmath,ivdoc}); and (iii) the public source code of the
scripts. Item (ii) is replaceable: any outward-rounded interval
implementation re-executing the same block lists obtains the same
verdicts. All scripts are plain Python 3.12 + \texttt{mpmath}
\cite{mpmath} (standard library \texttt{fractions} for the exact
algebra), run independently of one another.

\subsubsection{The inclusion property and tail bounds}\label{subsec:incl}

\begin{lemma}[Inclusion property, {\cite{MoKC}}]\label{lem:incl}
Let $E$ be any expression built from interval-enclosed constants, the
variable, interval arithmetic operations, and interval elementary
functions, all with outward-rounded semantics. Then for every
interval box $I$ in its domain, $E(I)$ contains the true value
$E(\tau)$ for every $\tau\in I$.
\end{lemma}
\begin{proof}
Induction on the structure of $E$; see \cite[Lemma~6.1]{Zf}.
\end{proof}

\begin{lemma}[Tail bound for eta-quotient products]\label{lem:tail}
Let $r:=\overline{|q|}<1$ be a certified upper bound. For
$d\ge1$, $e\in\Zz$ and the principal branch of $\log$,
\[
  \Bigl|\log\prod_{n>N}\bigl(1-q^{dn}\bigr)^{e}\Bigr|
  \le E:=\frac{|e|\,r^{d(N+1)}}{(1-r^d)(1-r)},
\]
so the truncated product carries the rigorous tail factor
$\exp(z)$ with $z$ enclosed in the box $\{u+\ii v:|u|,|v|\le E\}$
(evaluated as the interval exponential \emph{of that box}).
\end{lemma}
\begin{proof}
For $|z|\le r<1$, $|\log(1-z)|\le|z|/(1-r)$; sum the geometric
series, cf.\ \cite[Lemma~6.2]{Zf}. Products and quotients of such
factors (the parametrization \eqref{eq:cdef} is a quotient of two
products) are enclosed numerator- and denominator-wise before
division.
\end{proof}

On the certified paths $r\le\e^{-2\pi\cdot0.702}<0.0122$, and with
truncation $N=40$ the tail width satisfies $E<10^{-78}$, far below
the working precision.

\subsubsection{\texttt{cert2}: path certification}\label{subsec:cert2}

A \emph{path certificate} for a path $\gamma\subset\Hh$ is a finite
list of parameter blocks $I_\nu$ covering $\gamma$ (down to the
endpoint, if the endpoint lies in $V$), with complex intervals
$Z_\nu$ such that $c(I_\nu)\subset Z_\nu$ and every $Z_\nu$ is
certified disjoint from $\Ast$ --- i.e.\ the interval lower bound of
$|\ReT z|^{2/3}+|\ImT z|^{2/3}$ over $Z_\nu$ exceeds $4^{2/3}$,
a pure interval endpoint comparison. By Lemma~\ref{lem:incl} and
Lemma~\ref{lem:tail} this is a proof, not an estimate
(cf.\ \cite[Prop.~6.4]{Zf}).

\begin{remark}[Generation and termination of the adaptive search]
\label{rem:termin}
Blocks are found by adaptive bisection: accept a block when its
enclosure is certified outside $\Ast$, else bisect. Termination is
not part of the trusted base, but it is guaranteed: the natural
interval extension of an expression built from Lipschitz non-singular
operations on a box has width $O(\operatorname{diam})$ \cite{MoKC};
the factors of \eqref{eq:cdef} are non-singular ($|q|<1$,
$|1\pm q^{dn}|$ bounded away from $0$) and the tail width is
$<10^{-78}$ uniformly; and on each compact path piece the true image
has positive distance to $\Ast$. Once the enclosure width drops
below that margin, the block is accepted. (For a path ending at a
point of $c^{-1}(K)$ --- which does \emph{not} occur in this paper
--- the same machinery certifies down to a cap
$(y_0,y_0+\varepsilon_0]$, closed analytically by a first-order
Taylor expansion with a Cauchy-estimate remainder and a certified
enclosure of the derivative at the endpoint; see
\cite[\S4.3]{Zf}. Theorems~\ref{thm:B} and~\ref{thm:F} have their
endpoints strictly inside $V_4$, so no cap is needed.)
\end{remark}

\subsubsection{\texttt{cert0}: exact algebraic locks}\label{subsec:cert0}

The exactness of $s_4(\tau_0)=s$ follows a four-step pattern:
(i) \emph{integrality}: $s_4\in\Zz(\!(q)\!)$ (exact integer check),
so by Shimura's CM theorem \cite{Cox} the value of the eta-quotient
(a modular unit with integral $q$-expansion) at a CM point is an
algebraic integer in the corresponding ring class field; the class
number bounds the degree, and conjugates are located via the level
structure;
(ii) \emph{symmetric functions}: interval locks with rigorous tails
force the elementary symmetric functions of the conjugate values
(e.g.\ sum $S$ and product $P$ for a quadratic value) to be explicit
integers;
(iii) \emph{pinning}: an interval enclosure of the target with radius
smaller than half the separation of the candidate roots singles out a
unique root;
(iv) \emph{reality}: real $q$ at the chosen points
($q(\ii y)>0$, $q((1+\ii y)/2)<0$) makes the products real and
simplifies the locks.
No exact equality is ever inferred from a floating-point
approximation alone: each lock is a discreteness argument over an
explicit finite candidate set.

\subsubsection{\texttt{verify\_P1}: the three-track CM evaluation}
\label{subsec:verify}

The evaluation $\EK_4(\tau_0)=R$ (station S5) is certified by three
independent tracks, any one of which is sufficient in principle:
\begin{itemize}
\item \textbf{[X]/[S] exact algebra.} The lattice decomposition of
  $-T_1+4T_2$ into Hecke and Dirichlet sums, the inclusion--exclusion
  coefficients, the closed forms of even-character values
  $L(\chi,2)$ \cite{Wa}, the functional-equation constants, and the
  final coefficient comparison in a rational \emph{$e$-basis} are all
  carried out in rational/algebraic-integer arithmetic
  (\texttt{fractions.Fraction}); no interval estimates enter. The
  anchors of this track are quoted standard inputs listed explicitly
  in each application section (Hecke's theorem on gr\"ossencharacter
  theta series \cite{Miyake}, class numbers, Dirichlet functional
  equations, Shimura's CM theorem).
\item \textbf{[V] interval locks.} The same identity is re-proved by
  outward-rounded interval arithmetic: Poisson row sums with
  $\coth/\tanh$ closed rows and rigorous Fourier tails for the
  lattice sums, $E_1$-continued-fraction-bracketed Mellin integrals
  for the $L$-values, Euler--Maclaurin for the Dirichlet values.
  Achieved half-widths are quoted in each application section
  (down to $9.1\times10^{-53}$).
\item \textbf{mp cross-checks.} Independent $60$-digit
  (\texttt{mpmath}, non-interval) reimplementations of every
  quantity, agreeing typically to $10^{-60}$; these corroborate the
  implementation but do not enter the proof.
\end{itemize}
Newform identifications use Ligozat's eta-quotient criteria
\cite{Lig} (weight, character, cusp orders), exact-integer theta
identity checks to order $q^{60}$ against the Sturm bound
\cite{Sturm} (which is $\le6$ in all our cases), and the LMFDB
labels \cite{LMFDB}. Level and root number of each twist are
determined by the \emph{Fricke ratio}
$f\bigl(\ii/(\sqrt N\,y)\bigr)\big/\bigl(y^3 f(\ii y/\sqrt N)\bigr)$
evaluated at two ordinates --- see \S\ref{subsec:twisttrap} for why
the naive ``functional-equation residue'' test is vacuous.

\subsubsection{Script inventory}\label{subsec:scripts}

\begin{center}\footnotesize
\begin{tabular}{lp{9.6cm}}
\toprule
script & content \\
\midrule
\texttt{cert0\_s4\_s7pair.py} & exactness $s_4(\ii\sqrt7)$,
  $s_4(\ii\sqrt7/2)=8292456\pm3132675\sqrt7$ (Theorem A) \\
\texttt{verify\_P1\_n4\_s7pair.py} & Theorem A: exact separation
  track [S1]--[S10], FE proofs (\S\ref{subsec:app1FE}), interval
  locks [V0]--[V6] \\
\texttt{cert0\_s4\_m144.py} & exactness $s_4(\tau_1)=-144$
  (Theorem B) \\
\texttt{cert2\_path\_n4\_m144.py} & anchor box + certified path for
  Theorem B (template for Theorem F) \\
\texttt{verify\_P1\_n4\_m144.py} & Theorem B: 32 checks, four tracks
  including [G] Sturm-level newform identification and [V] interval
  locks \\
\texttt{n4\_m144\_true\_mahler.py} & true Mahler side of
  Theorem B by spectral torus integration (22-digit collision) \\
\texttt{cert0\_n4\_p3\_t1t2.py} & exactness of the $s_4$-values of
  Theorems C, D and of the conjugate of C (Theorem F's $s$-value) \\
\texttt{verify\_P1\_n4\_p3\_t1t2.py} & Theorems C, D: three tracks,
  [G] identifications, Fricke-ratio level determinations \\
\texttt{cert0\_n4\_p4\_t3.py} & Theorem E: four-point quartic lock
  of the $s_4$-value \\
\texttt{verify\_P1\_n4\_p4\_t3.py} & Theorem E: three tracks,
  genus-character decomposition, [V] locks \\
\texttt{n4\_p4\_t4\_cert.py} & certified path for Theorem F
  (interior-point verdict, margin $0.11590$) \\
\texttt{n4\_p4\_t4\_verify.py} & Theorem F: three tracks, shifted
  $\coth/\tanh$ row machine, [V] locks \\
\bottomrule
\end{tabular}
\end{center}

All computations use \texttt{mpmath} interval arithmetic at
$40$--$70$ decimal digits (\texttt{iv.dps}); the certification
scripts print \texttt{PASS}/\texttt{FAIL} per check and terminate
with an all-checks-passed line, re-verified by the author.

% ============================================================
\section{Application I: the \texorpdfstring{$\sqrt7$}{sqrt7} pair}
\label{sec:app1}

Both CM points are pure imaginary and deep inside Samart's region:
\[
  \tau_A=\ii\sqrt7\quad(\ImT\tau_A=\sqrt7>1/\sqrt2),\qquad
  \tau_B=\frac{\ii\sqrt7}{2}\quad(\ImT\tau_B=\tfrac{\sqrt7}{2}>1/\sqrt2),
\]
so Theorem~\ref{thm:samart4} applies directly: station S3 is free.
The work is S4 (the exact $s_4$-values) and S5 (the exact CM
evaluations, \emph{separately} for the two points --- an earlier
stage of this project had proved only the sum of the two identities;
the ray-class decomposition below closes each point independently).

\subsection{Station S4: the exact \texorpdfstring{$s_4$}{s4}-values}

By \texttt{cert0\_s4\_s7pair.py} (four-step lock of
\S\ref{subsec:cert0}): the two values are conjugate algebraic
integers in $\Qq(\sqrt7)$ (Shimura integrality; $q(\tau_A)>0$,
$q(\tau_B)>0$); interval locks force the symmetric functions to
integers,
\[
  S=16584912,\qquad P=69257922561=3^{12}\cdot19^{4}
  \quad\bigl(\text{exact check: }8292456^2-7\cdot3132675^2=P\bigr),
\]
and the root separation $1.6\times10^{7}$ dwarfs the lock radii
$|s_4(\ii\sqrt7)-r_+|\le9.6\times10^{-51}$,
$|s_4(\ii\sqrt7/2)-r_-|\le3.6\times10^{-54}$. Hence exactly
\[
  s_4(\ii\sqrt7)=8292456+3132675\sqrt7=:r_+,\qquad
  s_4(\ii\sqrt7/2)=8292456-3132675\sqrt7=:r_- .
\]

\subsection{Station S5: the lattice decomposition}

Write $\varpi=(1+\sqrt{-7})/2$, so $2\varpi=1+\sqrt{-7}$,
$K=\Qq(\sqrt{-7})$, $h(-7)=1$, $\OO_K^\times=\{\pm1\}$. The lattices
of \eqref{eq:EK4lattice} at the two points are (exact set equalities,
checks [S1]--[S2] of \texttt{verify\_P1\_n4\_s7pair.py}):
\begin{itemize}
\item at $\tau_A$: $\Lambda_1(\tau_A)=\Zz+\ii\sqrt7\Zz=\OO^{(2)}$
  (the order of conductor $2$, discriminant $-28$), and
  $\Lambda_2(\tau_A)=\OO^{(4)}$ (conductor $4$, discriminant $-112$);
\item at $\tau_B$: $\Lambda_2(\tau_B)=\OO^{(2)}$, and
  $\Lambda_1(\tau_B)=\tfrac12\Lambda'$, where
  $\Lambda'=2\Zz+\ii\sqrt7\Zz$ is the index-$2$ \emph{non-ideal}
  sublattice of $\OO^{(2)}$; by $(-4)$-homogeneity,
  $T_1(\tau_B)=16\,T(\Lambda')$.
\end{itemize}
With $B(\Lambda)=\sum'_{\lambda\in\Lambda}|\lambda|^{-4}$,
$G(\Lambda)=\sum'\bar\lambda^2|\lambda|^{-6}$ (real: all three
lattices are conjugation-stable), $T=2G+B$, the ray-class
decompositions
\begin{align*}
  \OO^{(2)}\setminus\{0\}&=\{\alpha\equiv1\ (2)\}\ \sqcup\
   2\OO_K\setminus\{0\},\\
  \OO^{(4)}&=2\OO^{(2)}\ \sqcup\ \{\alpha\equiv1,3\ (4)\},\\
  \Lambda'&=2\OO^{(2)}\ \sqcup\ \{\alpha\equiv1{+}2\varpi,\,
   3{+}2\varpi\ (4)\},
\end{align*}
the inclusion--exclusion coefficients $9/16$ and $23/16$ at the split
prime $(2)=(\varpi)(\bar\varpi)$ (exact, using
$\varpi^2+\bar\varpi^2=-3$), and the mod-$4$ ray-class pairing of the
character $\chi_{-4}\circ N$ (values $+1,+1,-1,-1$ on the four unit
classes, check [X4]) give, all in exact rational arithmetic
(checks [S3]--[S6]):
\[
\begin{array}{c|cc|c}
\text{lattice} & B & G & T=2G+B\\ \hline
\OO^{(2)} & \frac54\zeta_K(2) & 3L_3 & 6L_3+\frac54\zeta_K(2)\\
\OO^{(4)} & \frac{41}{64}\zeta_K(2)+z_V & \frac{13}{8}L_3+L_3^{\mathrm{tw}}
 & \frac{13}{4}L_3+2L_3^{\mathrm{tw}}+\frac{41}{64}\zeta_K(2)+z_V\\
\Lambda' & \frac{41}{64}\zeta_K(2)-z_V & \frac{13}{8}L_3-L_3^{\mathrm{tw}}
 & \frac{13}{4}L_3-2L_3^{\mathrm{tw}}+\frac{41}{64}\zeta_K(2)-z_V
\end{array}
\]
where $L_3:=L(g_7,3)$, $L_3^{\mathrm{tw}}:=L(g_7\otimes\chi_{-4},3)$,
and $z_V:=L(\chi_{-4},2)L(\chi_{-28},2)$. The pointwise assemblies
(check [S8], including the separation witness
$\mathrm{combA}\neq\mathrm{combB}$):
\begin{align*}
  \mathrm{combA}&:=-T\bigl(\OO^{(2)}\bigr)+4T\bigl(\OO^{(4)}\bigr)
   =7L_3+8L_3^{\mathrm{tw}}+\frac{21}{16}\zeta_K(2)+4z_V,\\
  \mathrm{combB}&:=-16\,T(\Lambda')+4T\bigl(\OO^{(2)}\bigr)
   =-28L_3+32L_3^{\mathrm{tw}}-\frac{21}{4}\zeta_K(2)+16z_V .
\end{align*}

\subsection{The functional equations and root numbers}\label{subsec:app1FE}

Two functional-equation inputs are proved in
\texttt{verify\_P1\_n4\_s7pair.py}, not quoted:
$M_7=\frac{7\sqrt7}{4\pi^3}L_3$ \cite[Lemma~2.8]{Zf} and the twisted
identity
\begin{equation}\label{eq:twistfricke}
  U_{a/4}\,W_{112}=4\,\gamma_a\,W_7\,U_{a/4}\qquad(a=1,3),
  \qquad
  \gamma_a=\begin{pmatrix}(1+7a^2)/4 & a\\ 7a & 4\end{pmatrix}
  \in\Gamma_0(7),
\end{equation}
an exact matrix identity (check [X10], with $\chi_{-7}(4)=+1$), which
gives $h|[W_{112}]_3=\ii\,h$ for $h=g_7\otimes\chi_{-4}$ --- the
Atkin--Li twist-$W$ technique of \cite{AL} made self-contained for
these parameters --- hence root number $+1$ at level $112$ and
$M_7^{\mathrm{tw}}=\frac{112\sqrt7}{\pi^3}L_3^{\mathrm{tw}}$. The
remaining quoted constants are the Dirichlet functional equations
(textbook) and the finite-sum evaluation
\[
  L(\chi_{-28},2)=\frac{2\pi^2}{7\sqrt7},
\]
whose character sum $\sum_a\chi_{-28}(a)a^2=448$ is an exact integer
check [X7] \cite{Wa}.

\subsection{Assembly}

With the $e$-basis
$e=\frac{\sqrt7}{\pi^3}\bigl(L_3,L_3^{\mathrm{tw}},\zeta_K(2),z_V\bigr)$
one has $M_7=\frac74e_1$, $M_7^{\mathrm{tw}}=112e_2$,
$d_7=\frac{21}{2}e_3$, $d_4=7e_4$ (exact, [S9]); the coefficient
comparison
\[
  \EK_4(\ii\sqrt7)=10\,\mathrm{combA}\cdot e
  =\Bigl(70,80,\tfrac{105}{8},40\Bigr)\cdot e
  =\frac{5}{28}\bigl(4M_7^{\mathrm{tw}}+224M_7+32d_4+7d_7\bigr),
\]
\[
  \EK_4\Bigl(\frac{\ii\sqrt7}2\Bigr)=5\,\mathrm{combB}\cdot e
  =\Bigl(-140,160,-\tfrac{105}{4},80\Bigr)\cdot e
  =\frac{5}{14}\bigl(4M_7^{\mathrm{tw}}-224M_7+32d_4-7d_7\bigr),
\]
is exact rational arithmetic (check [S10]), and the two right-hand
sides are distinct --- the separation is complete and each identity
holds individually. Combined with Theorem~\ref{thm:samart4} and
\S\ref{subsec:cert0} this proves Theorem~\ref{thm:A}. \qed

\begin{remark}[Independent confirmations]
The mp track (60 dps, independent implementations: Mellin
integrals, Poisson row sums, ray-class sums) confirms every step to
$9\times10^{-60}$; the interval-lock track [V0]--[V6]
(\texttt{iv.dps}$=70$: continued-fraction bracketing, Fourier tail
bounds, Euler--Maclaurin) locks the two final identities with
half-widths $1.15\times10^{-52}$ and $2.30\times10^{-52}$.
Reference values:
\begin{align*}
  M_7&=0.1026716077789020112104565948982929139989\ldots,\\
  M_7^{\mathrm{tw}}&=9.887687024790914246588215159482724883665\ldots
\end{align*}
\end{remark}

% ============================================================
\section{Application II: \texorpdfstring{$n_4(-144)$}{n4(-144)}}
\label{sec:app2}

The CM point
\[
  \tau_1=\frac{1+\sqrt{-3}}2
\]
is \emph{not} pure imaginary, so Theorem~\ref{thm:samart4} does not
apply directly; this is the first use of the full continuation
machine in the $n_4$ family. It also shows the machine in its
shortest form: the target is a strict interior point of $V_4$, so no
endpoint cap and no boundary-continuity step are needed.

\subsection{Station S3: the certified path}

The path $\gamma$ from the anchor box
$D=\{|\ReT\tau|\le1/32,\ |\ImT\tau-1|\le1/32\}$ to $\tau_1$ ---
horizontal at $\ImT\tau=1$ to $1/2+\ii$, then vertical at
$\ReT\tau=1/2$ down to $\tau_1$ --- is certified inside $V_4$ by
\texttt{cert2\_path\_n4\_m144.py} (machinery of
\S\ref{subsec:cert2}): the anchor box passes as a single block with
astroid margin $0.291$ and the certified bound $|s_4|\ge493>256$
(which activates Lemma~\ref{lem:anchor}: $F=G$ on $D$); the
horizontal leg takes $2$ blocks (minimum margin $0.148$); the
vertical leg takes $1$ block (margin $1.051$); and $\tau_1$ itself
satisfies $c(\tau_1)=2\sqrt3\,\e^{-\ii\pi/4}$, whose astroid
functional is $3.634>4^{2/3}=2.520$ --- a strict exterior point of
$\Ast$, so $\tau_1\in V_4$ itself and $\tau_1\in W$. The endpoint
Taylor cap of \cite{Zf} is absent. By Theorem~\ref{thm:main},
$n_4(s_4(\tau_1))=\EK_4(\tau_1)$.

\subsection{Station S4: the exact \texorpdfstring{$s_4$}{s4}-value}

\texttt{cert0\_s4\_m144.py}: $q(\tau_1)=-\e^{-\pi\sqrt3}$ is a
negative real, all products are real, and the value is pinned to a
\emph{rational integer} (Shimura integrality plus $h(-3)=1$, so the
ring class field is $K$; reality from the integral $q$-expansion and
$T$-invariance):
$|s_4(\tau_1)+144|\le2.3\times10^{-51}<1/2$, hence
$s_4(\tau_1)=-144$ exactly.

\subsection{Station S5: the CM evaluation}

$K=\Qq(\sqrt{-3})$, $\OO_K=\Zz[\omega]$,
$\omega=(-1+\sqrt{-3})/2$, units $\mu_6$. The structural
simplification is the \emph{$\mu_6$ annihilation}
$\sum_{u\in\mu_6}u^2=0$ (exact check [X1]): it forces
$G(\OO_K)=0$, which is also the reason $\Qq(\sqrt{-3})$ has no
conductor-$(1)$ gr\"ossencharacter of type $(2,0)$ and the evaluation
involves a non-trivial conductor. The lattices are
$\Lambda_1(\tau_1)=\OO_K$ (since $\tau_1=1+\omega=-\omega^2$) and
$\Lambda_2(\tau_1)=\OO_2$ (conductor $2$; $2$ is inert), giving
(exact [X2]--[X8]):
\[
  T_1(\tau_1)=6\zeta_K(2),\qquad
  T_2(\tau_1)=4L(g_{12},3)+\frac94\zeta_K(2),\qquad
  -T_1+4T_2=16\,L(g_{12},3)+3\zeta_K(2);
\]
note the $\zeta_K(2)$-term does \emph{not} cancel (coefficient $3$)
--- this is the source of the $d_3$-term. The theta identity
$2g_{12}=\sum'_{\alpha\equiv1\,(2)}\alpha^2q^{N\alpha}$ is checked
with exact integer arithmetic to order $q^{60}$ (check [L1]); the
newform identification
$g_{12}=\eta(2\tau)^3\eta(6\tau)^3\in S_3(\Gamma_0(12),\chi_{-3})$
is certified at Sturm level (checks [G1]--[G3]): Ligozat's criteria
\cite{Lig} (both congruences $24\equiv0\bmod24$; character
$-2^33^3$, square kernel $-3$), all six cusp orders equal $1$
(cuspidal; divisor degree $6=\frac{3}{12}\cdot24$), Sturm bound
$\lfloor3\cdot24/12\rfloor=6<60$ \cite{Sturm}; LMFDB label
\textsf{12.3.c.a} \cite{LMFDB}. With
$M_{12}=\frac{6\sqrt3}{\pi^3}L(g_{12},3)$ (root number $+1$, Fricke
ratio) and $d_3=\frac{3\sqrt3}{4\pi}L(\chi_{-3},2)$, the $e$-basis
assembly gives (exact, [S1]):
\[
  \EK_4(\tau_1)=\frac{5\sqrt3}{\pi^3}
   \bigl(16L(g_{12},3)+3\zeta(2)L(\chi_{-3},2)\bigr)
  =\frac{10}{3}\bigl(4M_{12}+d_3\bigr),
\]
both sides being $(80,\,5/2)$ in the basis
$\bigl(\sqrt3 L(g_{12},3)/\pi^3,\ \sqrt3 L(\chi_{-3},2)/\pi\bigr)$.
Theorem~\ref{thm:B} follows. \qed

\begin{remark}[Confirmations and the true-Mahler end]
The mp track confirms the identity to $3.7\times10^{-60}$ (check
[E1]); the interval-lock track [V0]--[V3] (shifted $\coth/\tanh$
rows, since $\ReT\tau_1=1/2$) locks it with half-width
$5.8\times10^{-52}$. Independently, direct spectral torus integration
(\texttt{n4\_m144\_true\_}\allowbreak\texttt{mahler.py}: $c(\tau_1)$ lies outside $\Ast$,
so the Jensen integrand is real-analytic periodic and the
two-dimensional trapezoidal rule converges spectrally; residual
$1.16\times10^{-22}$ at $N=128$) collides the two sides of
Theorem~\ref{thm:B} to $22$ digits:
\[
  n_4(-144)=5.09841965623737833323\ldots
\]
The propagation argument does not depend on this numerical value.
Reference value:
$M_{12}=0.3016149874129407464690529311477683998854\ldots$\,.
\end{remark}

% ============================================================
\section{Application III: class number one, deep interior points
(Theorems C and D)}\label{sec:app3}

Both CM points are pure imaginary and deep inside Samart's region,
\[
  \tau_C=\ii\sqrt2\quad(\ImT\tau_C=\sqrt2),\qquad
  \tau_D=2\ii\quad(\ImT\tau_D=2),
\]
so Theorem~\ref{thm:samart4} applies directly and the proofs consist
of \texttt{cert0} locks and three-track CM evaluations. The new
methodological content of this section is concentrated in station S5:
the character $\chi_8\circ N$ as an exact mod-$2$ ray-class projector,
and the level determination of the twists, which contains a genuine
trap.

\subsection{The ray-class projector \texorpdfstring{$\chi_8\circ N$}{chi8 o N}}
\label{subsec:rayproj}

In both fields the twist character materializes as a parity condition.
Exact checks ([X4], [Y4] of \texttt{verify\_P1\_n4\_p3\_t1t2.py}):
\begin{itemize}
\item in $K=\Qq(\sqrt{-2})$: for $a+b\sqrt{-2}$ of odd norm,
  $\chi_8\bigl(N(a+b\sqrt{-2})\bigr)=+1\iff b$ even ($a$ odd);
\item in $K=\Qq(\ii)$: for $a+b\ii$ of odd norm,
  $\chi_8\bigl(N(a+b\ii)\bigr)=+1\iff b\equiv0\ (4)$ ($a$ odd,
  $b$ even).
\end{itemize}
Hence $(1+\chi_8\circ N)/2$ is exactly the projector onto the mod-$2$
ray class $\{\alpha\equiv1\ (2)\}$, and in class number one no genus
character is needed: every ideal decomposition in this section is a
parity bookkeeping exercise certified in rational arithmetic. Also
used everywhere below: the closed form
\[
  L(\chi_8,2)=\frac{\pi^2\sqrt2}{16}
\]
(even-character finite-sum formula \cite{Wa}: $\tau(\chi_8)=2\sqrt2$,
$\sum_a\chi_8(a)a^2=16$, exact check [X7]).

\subsection{Theorem C: \texorpdfstring{$\tau_C=\ii\sqrt2$}{tauC}}
\label{subsec:app3C}

\textbf{Station S4.} \texttt{cert0\_n4\_p3\_t1t2.py}:
$q(\ii\sqrt2)=\e^{-2\pi\sqrt2}>0$ and
$q\bigl((1+\ii\sqrt2)/2\bigr)=-\e^{-\pi\sqrt2}<0$ are real, so real
interval locks suffice; the symmetric functions lock to the integers
\[
  S=7312,\qquad P=-153664=3656^2-2\cdot2600^2
\]
($|\Delta|\le1.1\times10^{-49}$, $5.2\times10^{-48}$; root separation
$7354$), pinning
\[
  s_4(\ii\sqrt2)=3656+2600\sqrt2,\qquad
  s_4\Bigl(\frac{1+\ii\sqrt2}2\Bigr)=3656-2600\sqrt2
\]
--- the second value is reused in Theorem~\ref{thm:F}.

\textbf{Station S5.} $K=\Qq(\sqrt{-2})$, $\OO_K=\Zz[\sqrt{-2}]$,
$h=1$, units $\pm1$, and $2=-(\sqrt{-2})^2$ ramifies. The lattices
are $\Lambda_1=\OO_K$ and $\Lambda_2=\OO_2=\Zz+2\OO_K$. With
$L_8:=L(g_8,3)$, $L_8^{\mathrm{tw}}:=L(g_8\otimes\chi_8,3)$, and
$z_V:=L(\chi_8,2)L(\chi_{-4},2)$, the inclusion--exclusion at the
ramified prime (factors $1/4$ for $B$, $-1/4$ for $G$) gives (exact,
[S]-track):
\[
  T(\OO_K)=4L_8+2\zeta_K(2),\qquad
  T(\OO_2)=\frac{11}{4}L_8+2L_8^{\mathrm{tw}}+\frac78\zeta_K(2)+z_V,
\]
\[
  \mathrm{comb}:=-T_1+4T_2
  =7L_8+8L_8^{\mathrm{tw}}+\frac32\zeta_K(2)+4z_V .
\]
The newform identification
$g_8=\eta(\tau)^2\eta(2\tau)\eta(4\tau)\eta(8\tau)^2
\in S_3(\Gamma_0(8),\chi_{-8})$ is Sturm-level
([G]-track: Ligozat congruences $24\equiv0\ (24)$; cusp orders
$(1,\tfrac12,\tfrac12,1)$; weighted divisor degree
$3=\frac{3}{12}\cdot12$; Sturm bound $3$; theta identity
$2g_8=\sum'_{\alpha\in\OO_K}\alpha^2q^{N\alpha}$ exact to $q^{60}$).
In the $e$-basis
$e=\frac{\sqrt2}{\pi^3}\bigl(L_8,L_8^{\mathrm{tw}},\zeta_K(2),z_V\bigr)$:
\[
  M_8=4e_1,\qquad M_8^{\mathrm{tw}}=32e_2\quad
  [\text{level }32,\ w=+1,\ \S\ref{subsec:twisttrap}],\qquad
  d_8=24e_3,\qquad d_4=16e_4,
\]
and the exact comparison ([S3])
$\EK_4(\ii\sqrt2)=10\,\mathrm{comb}\cdot e=(70,80,15,40)\cdot e$
is the right-hand side of Theorem~\ref{thm:C}. \qed

\subsection{Theorem D: \texorpdfstring{$\tau_D=2\ii$}{tauD}}
\label{subsec:app3D}

\textbf{Station S4.} The same script locks
\[
  S=286416,\qquad P=-126023688=143208^2-2\cdot101574^2
\]
($|\Delta|\le4.1\times10^{-48}$, $3.7\times10^{-45}$; separation
$2.87\times10^{5}$), pinning
$s_4(2\ii)=143208+101574\sqrt2$ and
$s_4\bigl((1+2\ii)/2\bigr)=143208-101574\sqrt2$ (the latter is the
open entry \#1 of Table~\ref{tab:open}).

\textbf{Station S5.} $K=\Qq(\ii)$, $\OO_K=\Zz[\ii]$, $h=1$, units
$\mu_4$ --- and $\sum_{u\in\mu_4}u^2=0$, so $G(\OO_K)=0$ (exact [Y1];
same annihilation mechanism as Theorem~\ref{thm:B}); $2=-\ii(1+\ii)^2$
ramifies. The lattices are $\Lambda_1=\OO_2$ and $\Lambda_2=\OO_4$.
Every ideal coprime to $2$ has exactly two generators
$\equiv1\ (2)$ ([Y3]), so the mod-$2$ ray sums carry \emph{no twist
term}: with $L_{16}:=L(g_{16},3)$,
$L_{16}^{\mathrm{tw}}:=L(g_{16}\otimes\chi_8,3)$, and
$z_{V2}:=L(\chi_8,2)L(\chi_{-8},2)$,
\[
  T(\OO_2)=4L_{16}+\frac74\zeta_K(2),\qquad
  T(\OO_4)=\frac94L_{16}+2L_{16}^{\mathrm{tw}}+\frac{55}{64}\zeta_K(2)
   +z_{V2},
\]
\[
  \mathrm{comb}=5L_{16}+8L_{16}^{\mathrm{tw}}+\frac{27}{16}\zeta_K(2)
   +4z_{V2},
\]
where the $\OO_4$ decomposition uses
$\OO_4\setminus\{0\}=\{a\text{ odd},\ b\equiv0\ (4)\}\sqcup
2\OO_2\setminus\{0\}$ and the projector of
\S\ref{subsec:rayproj} ([Y4]). The newform identification
$g_{16}=\eta(4\tau)^6\in S_3(\Gamma_0(16),\chi_{-4})$ is Sturm-level
(five cusp orders all $1$; weighted degree $6=\frac{3}{12}\cdot24$;
Sturm bound $6$; theta identity
$2g_{16}=\sum_{\alpha\equiv1\,(2)}\alpha^2q^{N\alpha}$ of conductor
$(2)$ exact to $q^{60}$). In the $e$-basis
$e=\frac1{\pi^3}\bigl(L_{16},L_{16}^{\mathrm{tw}},\zeta_K(2),z_{V2}\bigr)$:
\[
  M_{16}=16e_1,\qquad M_{16}^{\mathrm{tw}}=128e_2\quad
  [\text{level }64,\ w=+1],\qquad d_4=12e_3,\qquad d_8=64e_4,
\]
and the exact comparison ([S6])
$\EK_4(2\ii)=20\,\mathrm{comb}\cdot e=(100,160,135/4,80)\cdot e$
is the right-hand side of Theorem~\ref{thm:D}. \qed

\subsection{The twist-level trap and the Fricke ratio}\label{subsec:twisttrap}

Two methodological points emerged here and now belong to the machine.

\textbf{(i) The ``FE residue'' test is vacuous.} A widespread
heuristic decides the root number of a completed $L$-function
$\Lambda(f,s)=N^{s/2}(2\pi)^{-s}\Gamma(s)L(f,s)$ by splitting the
Mellin integral at $y=1$ and comparing the two sides of
$\Lambda(1)\stackrel?=w\Lambda(2)$. But the split gives
$\Lambda(s)=I_1(s)+wI_2(3-s)$-type expressions for which, at
$w=+1$, the two sides are \emph{the same linear combination of the
same two integrals at every candidate level $N$}: the ``residue''
vanishes identically in $N$ and certifies nothing. (The same vacuity
affected a check in the $\sqrt7$-pair script; there the conclusion
was independently secured by the Fricke ratio, so no result was
endangered.)

\textbf{(ii) The non-vacuous discriminator is the Fricke ratio.}
For the true level $N_0$ one has the exact transformation
$f\bigl(\ii/(\sqrt{N_0}\,y)\bigr)=w\,y^3f\bigl(\ii y/\sqrt{N_0}\bigr)$
for \emph{all} $y>0$, coming from the Fricke involution (proved for
each form in this paper either by the eta transformation law or by a
twisted Fricke matrix identity such as \eqref{eq:twistfricke}). Hence
the ratio
\[
  R_N(y):=\frac{f\bigl(\ii/(\sqrt N\,y)\bigr)}
                 {y^3f\bigl(\ii y/\sqrt N\bigr)}
\]
is identically $w$ at $N=N_0$, while at $N\neq N_0$ it varies with
$y$; evaluation at two ordinates separates the candidates. This is
how the levels in this paper were fixed: $32$ for $g_8\otimes\chi_8$,
$64$ for $g_{16}\otimes\chi_8$ (two-ordinate agreement to
$3.1\times10^{-61}$ and better), and $96$ for the twists of
Section~\ref{sec:app4} (unique match in the candidate scan
$\{24,48,96,192,384\}$).

\textbf{(iii) The trap itself.} For $g_{16}\otimes\chi_8$ the naive
level guess $\operatorname{lcm}(16,4^2)$-style reasoning suggests
$32$ --- and the vacuous test (i) happily ``confirms'' it. With the
wrong level-$32$ Mellin constant, $M_{16}^{\mathrm{tw}}$ comes out
smaller by a factor $2\sqrt2$, and the conjectured identity fails
numerically by $\sim1\%$; a first pass over Theorem~\ref{thm:D}
recorded exactly this mismatch, briefly suggesting an $s=81$-type
refutation. The Fricke ratio at two ordinates gives the true level
$64$, and with it the identity holds to $2.5\times10^{-60}$. Since a
wrong level \emph{silently scales} an $L$-value by a simple algebraic
factor, we record the rule: \emph{never determine a twist level by a
residual; determine it by the Fricke ratio at two ordinates, and
cross-check against the assembled identity.}

\begin{remark}[Confirmations for Theorems C and D]
The mp tracks (60 dps) confirm the two final identities to
$2.5\times10^{-60}$ each; the interval-lock tracks
[V] contain zero with half-widths $\le4.6\times10^{-53}$ (Theorem C)
and $\le6.2\times10^{-53}$ (Theorem D). Direct torus integration
($c\approx9.25$ and $23.14$, smooth integrands) confirms both at
float64 level. Reference values:
\begin{align*}
  M_8&=0.1353184954231269106150060654595395883532\ldots,\\
  M_8^{\mathrm{tw}}&=1.549084847095611486271658159444376281785\ldots
   \quad(\text{level }32,\ w=+1),\\
  M_{16}&=0.4975478545679851929276297713583412482363\ldots,\\
  M_{16}^{\mathrm{tw}}&=4.336361249244445609135012019721278212973\ldots
   \quad(\text{level }64,\ w=+1),\\
  d_8&=\frac{4\sqrt2}{\pi}L(\chi_{-8},2)
   =1.9171950931209540617988237536697845644585\ldots
\end{align*}
($d_8$ by the functional equation and by direct numerical
differentiation, agreeing exactly at working precision).
\end{remark}

% ============================================================
\section{Application IV: class number two (Theorem E)}\label{sec:app4}

The CM point $\tau_E=\ii\sqrt6$ is pure imaginary with
$\ImT\tau_E=\sqrt6>1/\sqrt2$, so Theorem~\ref{thm:samart4} applies
directly. The field $K=\Qq(\sqrt{-6})$ has discriminant $-24$ and
\emph{class number $2$} --- the first such case treated by the
machine --- and everything is one step richer: two newforms, a genus
character, a non-principal ideal class, and a four-point
\texttt{cert0} lock.

\subsection{Station S4: the four-point quartic lock}

Class number two makes the $s_4$-value quartic over $\Qq$, in
$\Qq(\sqrt2,\sqrt3)$. \texttt{cert0\_n4\_p4\_t3.py} locates and
locks all four conjugate values:
\begin{center}
\begin{tabular}{lll}
\toprule
signs & CM point & lock half-width\\
\midrule
$(+,+)$ & $\ii\sqrt6$ \ ($q=\e^{-2\pi\sqrt6}>0$) & $7.0\times10^{-47}$\\
$(-,+)$ & $\ii\sqrt6/4$ \ ($q=\e^{-\pi\sqrt6/2}>0$) & $4.4\times10^{-51}$\\
$(+,-)$ & $(1+\ii\sqrt6)/2$ \ ($q=-\e^{-\pi\sqrt6}<0$) & $3.1\times10^{-50}$\\
$(-,-)$ & $1/2+\ii\sqrt6/4$ \ ($q=-\e^{-\pi\sqrt6/2}<0$) & $6.8\times10^{-53}$\\
\bottomrule
\end{tabular}
\end{center}
with minimal polynomial
\[
  X^4-4829472X^3-8676927360X^2+3042408646656X+7520406736896,
\]
whose coefficients are confirmed as integers both by exact
$\Zz[\sqrt2,\sqrt3]$ arithmetic and by the interval locks; the
minimum root separation $305.0$ dwarfs every lock radius. Two
structural observations aided the location (used for location only,
not in the certified chain): $s_4$ is $W_2$-invariant,
$s_4(-1/(2\tau))=s_4(\tau)$, observed numerically to $10^{-33}$ at
sample points; and the second class of discriminant $-24$ degenerates:
$s_4(\ii\sqrt6/2)=s_4(\ii\sqrt6/6)=2304\in\Zz$, which is why that
value is absent from the quartic roots. In particular the $(+,-)$
value is open entry \#2 of Table~\ref{tab:open} and the $(-,\pm)$
values are entries \#3, \#4: their station S4 is already done.

\subsection{The two newforms and the genus character}\label{subsec:genus}

The class group is $\Zz/2$ and coincides with the genus group. The
theta series of the two gr\"ossencharacters $\psi$ and
$\psi\cdot\chi_g$ of type $(2,0)$ have coefficients (exact-integer
checks [th0]--[th6]: integrality, $a(1)=1$, the Hecke recursion
$a(p^2)=a(p)^2-\chi_{-24}(p)p^2$, multiplicativity, vanishing at
inert primes, distinctness):
\[
  a^{(1)}(n)=P(n)-Q(n),\qquad a^{(2)}(n)=P(n)+Q(n),
\]
\[
  P(n)=\frac12{\sum_{a^2+6b^2=n}}^{\!\prime}\bigl(a^2-6b^2\bigr),
  \qquad
  Q(n)=\frac14{\sum_{2x^2+3y^2=n}}^{\!\prime}\bigl(4x^2-6y^2\bigr),
\]
$P$ collecting the principal ideal class and $Q$ the non-principal
one (elements of norm $2n$ in $\mathfrak p_2=(2,\sqrt{-6})$ are in
bijection with non-principal ideals of norm $n$). By Hecke's theorem
\cite{Miyake} both are newforms of $S_3(\Gamma_0(24),\chi_{-24})$,
which has dimension $2$, so they exhaust the space; the coefficient
match ($q\mp2q^2\pm3q^3+4q^4\pm2q^5-6q^6$) fixes the LMFDB labels
\[
  g_{24}^{(1)}=P-Q=\textsf{24.3.h.a},\qquad
  g_{24}^{(2)}=P+Q=\textsf{24.3.h.b},
\]
the unique rational CM newforms of level $24$ and weight $3$
\cite{LMFDB} (in agreement with Sch\"utt's classification \cite{Sch}).
Samart's $M_{24}^{(1)}$ (coefficient $28$ in
Theorem~\ref{thm:E}) is our $g_{24}^{(1)}=P-Q$; swapping the labels
breaks the identity at the $0.4$ level, a labeling pitfall we record.
The \emph{genus character} is
$\chi_g=\chi_8\circ N=\chi_{-3}\circ N$ on ideal norms (exact check
[th2]: $\chi_8\chi_{-3}=\chi_{-24}$, the field character, so the two
compositions coincide on norms), and it handles exactly the
non-principal class that the mod-$2$ ray projector of
\S\ref{subsec:rayproj} cannot see.

\subsection{Station S5: the class-number-two lattice decomposition}
\label{subsec:class2}

The lattices are $\Lambda_1=\OO_K$ and $\Lambda_2=\OO_2$. The new
ingredients, relative to class number one, are: (a) the class-group
projector $(1+\chi_g)/2$ on the $B$-side and the difference of the
two $\psi$-values on the $G$-side; (b) sums over $\mathfrak p_2$ as a
lattice, with homothety factors $1/4$ (not $1/2$ --- numerically
locked, then certified); (c) the ray condition
$\{a\text{ odd},b\text{ even}\}$ via $\chi_{-8}\circ N$, which is
\emph{not} a Hecke character (nontrivial on principal ideals) but
whose $L_K$-value still factors as a product of Dirichlet
$L$-functions. Exact decomposition:
\begin{align*}
  B(\OO_K)&=\zeta_K(2)+L(\chi_8,2)L(\chi_{-3},2),&
  G(\OO_K)&=L\bigl(g_{24}^{(1)},3\bigr)+L\bigl(g_{24}^{(2)},3\bigr),\\
  B(\mathfrak p_2)&=\frac{\zeta_K(2)-L(\chi_8,2)L(\chi_{-3},2)}4,&
  G(\mathfrak p_2)&=\frac{L(g_{24}^{(2)},3)-L(g_{24}^{(1)},3)}4,
\end{align*}
(the sign in $G(\mathfrak p_2)$ is $\psi(\mathfrak p_2)=-2$), and
\[
  B_{\mathrm{odd},\chi}=L(\chi_{-8},2)L(\chi_{3},2)
   +L(\chi_{-4},2)L(\chi_{24},2),\qquad
  G_{\mathrm{odd},\chi}=L\bigl(g_{24}^{(1)}\otimes\chi_{-8},3\bigr)
   +L\bigl(g_{24}^{(2)}\otimes\chi_{-8},3\bigr).
\]
Note $T(\Lambda)=B(\Lambda)+2G(\Lambda)$ for \emph{every} lattice:
the pointwise identity $4(\ReT z)^2/|z|^6-1/|z|^4
=2\ReT(z^2)/|z|^6+1/|z|^4$ (exact check [X0]) plus absolute
convergence, closing what had briefly appeared to be a gap. The even
primitive closed forms \cite{Wa}
\[
  L(\chi_3,2)=\frac{\pi^2\sqrt3}{18}\quad
  (\text{conductor }12,\ \textstyle\sum a^2\chi_3(a)=48),\qquad
  L(\chi_{24},2)=\frac{\pi^2\sqrt6}{24}\quad
  (\textstyle\sum a^2\chi_{24}(a)=288)
\]
are absorbed by the $d_8$ and $d_4$ terms in the assembly. With
$L_1,L_2,L_1^{\mathrm{tw}},L_2^{\mathrm{tw}}$ the four $L$-values at
$s=3$,
\begin{align*}
  \mathrm{comb}=\;&\frac72L_1+\frac32L_2+4L_1^{\mathrm{tw}}
   +4L_2^{\mathrm{tw}}+\frac34\zeta_K(2)\\
   &+\frac74L(\chi_8,2)L(\chi_{-3},2)
   +2L(\chi_{-8},2)L(\chi_3,2)+2\,\mathrm{Cat}\,L(\chi_{24},2),
\end{align*}
where $\mathrm{Cat}=L(\chi_{-4},2)$ is Catalan's constant,
and the [S1]--[S8] Fraction-exact assembly (single-term substitutions)
produces the coefficient list
\[
  \tfrac{35}{12},\ \tfrac54,\ \tfrac{5}{12},\ \tfrac{5}{12},\
  \tfrac{35}{12},\ \tfrac52,\ \tfrac56,\ \tfrac{5}{48},
\]
yielding Theorem~\ref{thm:E}. \qed

\begin{remark}[Traps recorded for the machine]
(i) $\chi_3$ has conductor $12$ with period values $0$ at even
positions; a wrongly tabulated period (nonzero even positions) makes
standard Dirichlet-$L$ software silently return a wrong value
($1.1325$ instead of $0.951\ldots$), which once produced a false
failure at the $3.5\times10^{-2}$ level. The machine now treats
character tables as exact data to be certified ([X3]). (ii) The
twist levels ($96$) were fixed by the Fricke-ratio scan of
\S\ref{subsec:twisttrap}, never by residues.
\end{remark}

\begin{remark}[Confirmations]
The mp track (60 dps) confirms the final identity to
$2.49\times10^{-60}$ (check [E]); the [V] interval-lock track
(Dual-number $\coth$ rows, $E_1$-bracketed Mellin with the
self-contained coefficient bound $|a(n)|\le6n^3$, level-$96$ Mellin
truncation $n_0\approx481$) locks $T(\OO_K)$, $T(\OO_2)$, and the
final identity with half-widths $4.50\times10^{-53}$,
$4.92\times10^{-53}$, and $9.72\times10^{-53}$. Direct torus
integration ($c\approx46.88$, smooth) confirms at float64 level.
Reference values (50 dps, Mellin; level $96$, $w=+1$ for the twists):
\begin{align*}
  M_{24}^{(1)}&=0.8366946514850983661417153108277990376289\ldots,\\
  M_{24}^{(2)}&=1.093580625535742995623816182417982487423\ldots,\\
  M_{24}^{(1),\mathrm{tw}}&=8.547412838595950972260990976022241348897
   \ldots,\\
  M_{24}^{(2),\mathrm{tw}}&=7.183569611977336512537940485577907159979
   \ldots
\end{align*}
\end{remark}

% ============================================================
\section{Application V: the boundary point that turned out interior
(Theorem F)}\label{sec:app5}

The point
\[
  \tau_4=\frac{1+\sqrt{-2}}2,\qquad \ImT\tau_4=\frac1{\sqrt2},
\]
lies \emph{exactly} on the boundary of Samart's applicability region.
Moreover $q(\tau_4)=-\e^{-\pi\sqrt2}<0$ is a negative real (outside
his $q\in(0,1)$ hypothesis) and $|s_4(\tau_4)|=20.96\ldots<256$
(outside the ${}_5F_4$ convergence region), so
Theorem~\ref{thm:samart4} fails to cover $\tau_4$ for two independent
reasons. This was expected to require a boundary-continuity argument.
It does not: the boundary in question is not the boundary of the
problem.

\subsection{Samart's boundary is not the topological boundary of
\texorpdfstring{$V_4$}{V4}}\label{subsec:notboundary}

A $60$-digit scan along the line $\ReT\tau=1/2$ shows the astroid
functional $|\ReT c|^{2/3}+|\ImT c|^{2/3}$ of $c(\tau)$ decreasing
monotonically from $4.377$ at $y=1$, crossing the threshold
$4^{2/3}=2.51984$ only at $y\approx0.693$ (margin $+0.064$ at
$y=0.700$, $-0.011$ at $y=0.690$). At $\tau_4$ ($y=1/\sqrt2=0.70711$)
the functional is $2.6357$, with margin $+0.116$. Since $c$ is
holomorphic (an open map) and $\Ast$ is closed, $\tau_4$ is a
\emph{strict interior point} of $V_4$. Samart's line
$\{\ImT\tau=1/\sqrt2\}$ is dictated by the real-$q$ mechanism of his
proof; the topological boundary of $V_4$ along this segment lies
$\approx0.014$ lower, and $\tau_4$ sits strictly between the two.
Consequently the short \texttt{cert2} path of Theorem~\ref{thm:B}
suffices, and no boundary-continuity step is needed.

\subsection{Station S3: the certified path}

\texttt{n4\_p4\_t4\_cert.py} (template
\texttt{cert2\_path\_n4\_m144.py}; leg B extended down to
$y=0.702<1/\sqrt2$ so that $\tau_4$ lies strictly inside the
certified leg; the $|q|$ bound $0.0122$ keeps the tail
$r^{41}<10^{-78}$):
\begin{itemize}
\item anchor box $D$: margin $0.29143$, certified $|c|\ge4.71$,
  $|s_4|\ge493.6>256$ (activating Lemma~\ref{lem:anchor});
\item leg A ($y=1$, $x\in[0,1/2]$): $2$ blocks, minimum margin
  $0.14759$;
\item leg B ($x=1/2$, $y\in[0.702,1]$): $4$ blocks, minimum margin
  $0.02270$ (most dangerous block $y\in[0.702,0.7393]$);
\item point enclosure at $\tau_4$: astroid margin $0.11590>0$ ---
  an interval-certified interior-point verdict;
\item self-tests: $21/21$ point inclusions (80-dps values inside
  interval enclosures, including $\tau_4$ itself) and $5/5$
  enclosure-containment checks (point enclosures inside block
  enclosures).
\end{itemize}
Hence the whole path and $\tau_4$ lie in $W$, the component of the
anchor box, and Theorem~\ref{thm:main} gives
$n_4(s_4(\tau_4))=\EK_4(\tau_4)$. As a byproduct, the open entries
\#1 ($\tau=(1+2\ii)/2$, the top corner of the legs) and \#7
($\tau=(3+\ii\sqrt{21})/6=1/2+0.7638\ii$, on leg B) of
Table~\ref{tab:open} already have their station S3 certified by this
same path.

\subsection{Station S4}

The value $s_4(\tau_4)=3656-2600\sqrt2$ is already locked exactly by
\texttt{cert0\_n4\_p3\_t1t2.py} as the conjugate point of
Theorem~\ref{thm:C} ($S=7312$, $P=-153664$;
\S\ref{subsec:app3C}).

\subsection{Station S5: a different lattice decomposition, not a
change of clothes}

Although the $s$-value is conjugate to Theorem~\ref{thm:C}'s, the CM
point is different, and the lattice structure differs:
$2\tau_4=1+\sqrt{-2}$ gives $\Lambda_2=\OO_K$ (not $\OO_2$), and
\[
  2\Lambda_1=2\Zz+\Zz(1+\sqrt{-2})
  =\{a+b\sqrt{-2}:a\equiv b\ (2)\}=:\OO'
  =2\OO_K\ \sqcup\ \bigl((1+\sqrt{-2})+2\OO_K\bigr),
\]
a non-ideal sublattice (exact [X1]--[X3]). The parity-class sums from
Theorem~\ref{thm:C}'s anchors ([X4], [X5], homothety $1/16$ [X8]):
\[
  B_{00}=\frac{\zeta_K(2)}8,\quad B_{11}=\frac34\zeta_K(2)-z_V,\qquad
  G_{00}=\frac{L_8}8,\quad G_{11}=\frac54L_8-L_8^{\mathrm{tw}},
\]
\[
  T(\OO')=\frac{11}{4}L_8-2L_8^{\mathrm{tw}}+\frac78\zeta_K(2)-z_V,
  \qquad T(\Lambda_1)=16\,T(\OO'),
\]
\[
  \mathrm{comb}=-T(\Lambda_1)+4T(\Lambda_2)
  =-28L_8+32L_8^{\mathrm{tw}}-6\zeta_K(2)+16z_V .
\]
Note the sign pattern: the twisted and $\zeta_K(2)$ terms flip sign
relative to Theorem~\ref{thm:C} (combinations $-28/+32$ versus
$+7/+8$), and the $-d_8$ term traces to the coefficient $-16$ of
$z_V$ (exact witness [S7]: the two sign patterns are genuinely
different). The $e$-basis assembly of \S\ref{subsec:app3C} gives
(exact [X10]/[S6])
\[
  \EK_4(\tau_4)=5\,\mathrm{comb}\cdot e=(-140,160,-30,80)\cdot e
  =\frac54\bigl(4M_8^{\mathrm{tw}}-28M_8+4d_4-d_8\bigr),
\]
proving Theorem~\ref{thm:F}. \qed

\begin{remark}[Confirmations]
The mp track (60 dps; the $L$-values recomputed in-script at levels
$8/32$, cross-checked against Section~\ref{sec:app3} to
$1.6\times10^{-40}$) confirms the final identity to
$9.33\times10^{-61}$ (check [E1]); the [V] interval-lock track ---
extended by a shifted $\coth/\tanh$ row machine
(\texttt{lattice\_T\_iv\_t4}: at $x_0=1/2$ the Poisson kernel rows
alternate between $\coth(\pi y)$ and $\tanh(\pi y)$; the power terms
are shift-independent and the tail bounds uniform in the shift) ---
locks $T(\Lambda_2)$, $T(\Lambda_1)$, and the final identity with
half-widths $3.48\times10^{-53}$, $6.94\times10^{-52}$, and
$9.12\times10^{-53}$. Direct torus integration at the complex
parameter $c=1.5128932208(1-\ii)$ (outside $\Ast$: functional
$2.6357$) confirms at float64 level:
$n_4(3656-2600\sqrt2)\approx3.5283920696\approx\EK_4(\tau_4)$
(the two computations differ by $2.5\times10^{-13}$, the float64
noise floor).
\end{remark}

% ============================================================
\section{The boundary of the method}\label{sec:boundary}

We collect the precise limits of the machine, all visible already in
Table~\ref{tab:open}.

\subsection{Two-dimensional critical images}

The decisive invariant is the dimension of the critical image
(Remark~\ref{rem:2D}). In the $n_2$ family the image is the slit
$[0,64]$, which does not separate $\Cc$, and every conjectured entry
was reachable: the family is complete \cite{Zf,ZGQ,GPQ}. In the
$n_4$ family the image is the astroid disc \eqref{eq:astroid}, which
has interior. A CM parameter with $c(\tau_0)$ in the interior of
$\Ast$ lies inside an open bad region of the $\tau$-plane (open
mapping), no certified path can approach it, and
Theorem~\ref{thm:main} is silent. The extreme case is the refuted
conjecture $n_4(81)=40M_7$ \cite{Zf2}: $c=3$ is an interior lattice
point of the astroid, the conjectured value equals the exact EK
evaluation at the attached CM point on a \emph{wrong sheet}, and
direct torus integration gives the different value
$n_4(81)=4.1655349907533676508(5)$. Entries \#4 and \#5 of
Table~\ref{tab:open} have $c(\tau_0)$ inside $\Ast$ ($s\approx-2.45$
and $s\approx-6.17$ respectively) and are blocked by the same
mechanism; entry \#5 is moreover exterior (see the next subsection).

\subsection{The wrong sheet below \texorpdfstring{$\ImT\tau=1/\sqrt2$}{Im tau = 1/sqrt2}}

Even where the critical image is avoided, the EK series itself can
fail. A systematic scan (\texttt{diag\_n4\_}\allowbreak\texttt{astroid.py}) shows
$\EK_4(\tau)\neq4\m(P+c(\tau))$ \emph{at every tested point} below
the line $\ImT\tau=1/\sqrt2$ --- including points with $c(\tau)$ real and
larger than $4$, where the ${}_5F_4$ evaluation converges and the
Mahler side is perfectly smooth (e.g.\ a difference of $11.4$ at
$\tau=0.3\ii$). The difference grows continuously from $0$ at the
line; there is no jump. The mechanism is the same wrong-sheet
phenomenon analyzed for the $n_2$ family in \cite[Remark~5.6]{Zf}:
the single-valued series $\EK_4$ follows a companion branch of
$\mt$ once the branch locus is crossed. The propagation of
Theorem~\ref{thm:main} cannot cross, because the identity theorem
fixes $\Psi\equiv0$ only on the anchor's component $W$: along the
imaginary axis the point $\tau=\ii/\sqrt2$ (where $c=4$, a cusp of
$\Ast$) separates $W$ from the lower component $W'$, and on $W'$ the
difference $H=F-G$ is non-constant. Entries \#3, \#9, and \#10 of
Table~\ref{tab:open} (pure imaginary, $y=0.612$, $0.154$, and
$0.463$, with $c>4$ real in each case) are blocked at this
\emph{series} level: the conjectured
identity --- a statement about the true Mahler measure, which
Samart's PSLQ evidence supports --- has no valid series input for
station S5, and proving it requires a corrected expression, e.g.\ a
modified Mahler measure of Samart--Tao $\widetilde n$-type \cite{ST},
whose design principle (choose the linear combination so that the
non-closed Deninger path \cite{De} becomes a meaningful pairing) is
compatible with stations S2/S4 of the machine.

\subsection{Open directions}

\begin{itemize}
\item \emph{The ready entries} \#1, \#2, \#6, \#7, \#8 of
  Table~\ref{tab:open}. For \#1 and \#7 the continuation step is
  already certified (Section~\ref{sec:app5}); what remains is
  mechanical: \texttt{cert0} pair locks in discriminants $-15$ and
  $-21$ (class number $2$; the genus machinery of
  Section~\ref{sec:app4} applies) and three-track P1 assemblies at
  levels $15$, $84$, and $24$, $16$.
\item \emph{The blocked entries} \#3, \#4, \#5, \#9, \#10: a
  corrected (wrong-sheet-free) formula is needed before the machine
  has anything to certify; the $\widetilde n$ framework of
  \cite{ST} is the natural candidate.
\item \emph{Samart's three-modular-$L$-value hypothesis}
  $s=16+1600\sqrt[3]2-1280\sqrt[3]4$ \cite{Sa15}: the conjectured
  value involves three modular $L$-values and no attached CM point in
  the $s_4$-image is known; the machine's station S1 has no input.
\item \emph{The Boyd-lineage entries}: the two-variable
  conjectures of Boyd's tables (e.g.\ conductor-$14$ and related
  families) have no known CM evaluation, so the machine does not
  apply; regulator-based approaches \cite{De,BZ}, and in particular
  deconditioning of the currently conditional identities in the
  Brunault--Zudilin framework \cite{BZ}, remain the promising route.
\end{itemize}

% ============================================================
\section*{Statements and Declarations}
\paragraph{Funding.} The author did not receive support from any
organization for the submitted work.
\paragraph{Competing Interests.} The author has no relevant financial or
non-financial interests to disclose.
\paragraph{Data Availability.} All certification scripts, frozen
certificates, and run logs are available at
\url{https://github.com/huiminZheng-collab/samart-mahler}, archived at
\url{https://doi.org/10.5281/zenodo.21711884}.
\paragraph{Use of AI tools.} Documented in the Declaration on the use of
AI tools on the title page.

\begin{thebibliography}{99}

\bibitem{Sa13}
D.~Samart,
\emph{Three-variable Mahler measures and special values of modular and
Dirichlet $L$-series},
Ramanujan J. \textbf{32} (2013), no.~2, 245--268.
arXiv:1205.4803. (Numbering of displayed results refers to the arXiv
version.)

\bibitem{Sa15}
D.~Samart,
\emph{Mahler measures as linear combinations of $L$-values of multiple
modular forms},
Canad. J. Math. \textbf{67} (2015), no.~2, 424--449.
arXiv:1303.6376.

\bibitem{HY}
Q.~He and D.~Ye,
\emph{On conjectures of Samart},
Manuscripta Math. \textbf{167} (2022), no.~3--4, 545--588.

\bibitem{Fei}
J.~Fei,
\emph{Mahler measure of 3D Landau--Ginzburg potentials},
Forum Math. \textbf{33} (2021), no.~5, 1369--1401.
arXiv:2009.09701.

\bibitem{Zf}
H.~Zheng,
\emph{Mahler measures at interior CM points: proofs of two conjectures
of Samart},
preprint, 2026; arXiv:2608.02255.

\bibitem{Zf2}
H.~Zheng,
\emph{Samart's conjecture $n_4(81)=40M_7$: the exact CM evaluation and
the two obstructions. A status report},
preprint, 2026; arXiv:2608.02265.

\bibitem{ST}
D.~Samart and Z.~Tao,
\emph{Determinants of Mahler measures and special values of
$L$-functions},
preprint; arXiv:2409.14714.

\bibitem{Ro}
M.~D.~Rogers,
\emph{New ${}_5F_4$ hypergeometric transformations, three-variable
Mahler measures, and formulas for $1/\pi$},
Ramanujan J. \textbf{18} (2009), 327--340.

\bibitem{RV}
F.~Rodriguez Villegas,
\emph{Modular Mahler measures, I},
in: Topics in Number Theory (University Park, PA, 1997),
Math.\ Appl.\ \textbf{467}, Kluwer, Dordrecht, 1999, 17--48.

\bibitem{ZGQ}
H.~Zheng, X.~Guo, and H.~Qin,
\emph{The Mahler measure of $(x+1/x)(y+1/y)(z+1/z)+\sqrt{k}$},
Electron. Res. Arch. \textbf{28} (2020), no.~1, 103--125.

\bibitem{GPQ}
X.~Guo, Y.~Peng, and H.~Qin,
\emph{Three-variable Mahler measures and special values of $L$-functions
of modular forms},
Ramanujan J. \textbf{54} (2021), no.~1, 147--175.

\bibitem{Cox}
D.~A.~Cox,
\emph{Primes of the Form $x^2+ny^2$: Fermat, Class Field Theory, and
Complex Multiplication}, 2nd ed., Wiley, 2013.

\bibitem{Miyake}
T.~Miyake,
\emph{Modular Forms},
Springer, Berlin, 1989.
(Theta series of Hecke gr\"ossencharacters, \S4.8.)

\bibitem{Lig}
G.~Ligozat,
\emph{Courbes modulaires de genre $1$},
Bull. Soc. Math. France, M\'em.\ \textbf{43}, 1975.
(Eta-quotient modularity and cusp-order criteria.)

\bibitem{Sturm}
J.~Sturm,
\emph{On the congruence of modular forms},
in: Number Theory (New York, 1984--1985),
Lecture Notes in Math.\ \textbf{1240}, Springer, Berlin, 1987,
275--280.

\bibitem{Wa}
L.~C.~Washington,
\emph{Introduction to Cyclotomic Fields}, 2nd ed.,
Grad.\ Texts in Math.\ \textbf{83}, Springer, New York, 1997.
(Finite-sum evaluation of $L(\chi,2)$ for even primitive $\chi$.)

\bibitem{AL}
A.~O.~L.~Atkin and W.~C.~W.~Li,
\emph{Twists of newforms and pseudo-eigenvalues of $W$-operators},
Invent. Math. \textbf{48} (1978), 221--243.

\bibitem{De}
C.~Deninger,
\emph{Deligne periods of mixed motives, $K$-theory and the entropy of
certain $\Zz^n$-actions},
J. Amer. Math. Soc. \textbf{10} (1997), 259--281.

\bibitem{BZ}
F.~Brunault and W.~Zudilin,
\emph{Many Variations of Mahler Measures: A Last Symphony},
Cambridge Univ. Press, Cambridge, 2020.

\bibitem{Sch}
M.~Sch\"utt,
\emph{CM newforms with rational coefficients},
Ramanujan J. \textbf{19} (2009), 187--205.

\bibitem{LMFDB}
The LMFDB Collaboration,
\emph{The $L$-functions and Modular Forms Database},
newforms \textsf{7.3.b.a}, \textsf{12.3.c.a}, \textsf{24.3.h.a},
\textsf{24.3.h.b}, \url{https://www.lmfdb.org}
(accessed August 2026).

\bibitem{mpmath}
F.~Johansson et al.,
\emph{mpmath: a Python library for arbitrary-precision floating-point
arithmetic} (version 1.3.0), \url{https://mpmath.org}
(accessed August 2026).

\bibitem{ivdoc}
The \texttt{mpmath} documentation, version 1.3.0,
\emph{Interval arithmetic (\texttt{mpmath.iv})},
\url{https://mpmath.org/doc/1.3.0/contexts.html}
(outward rounding of interval endpoints; accessed August 2026).

\bibitem{MoKC}
R.~E. Moore, R.~B. Kearfott, and M.~J.~Cloud,
\emph{Introduction to Interval Analysis},
SIAM, Philadelphia, 2009.

\end{thebibliography}

\end{document}
